% concatenated from main.tex, math_commands.tex
\documentclass{article}
\usepackage{mystyle,times}

%%%%% NEW MATH DEFINITIONS %%%%%

\usepackage{amsmath,amsfonts,bm}

% Mark sections of captions for referring to divisions of figures
\newcommand{\figleft}{{\em (Left)}}
\newcommand{\figcenter}{{\em (Center)}}
\newcommand{\figright}{{\em (Right)}}
\newcommand{\figtop}{{\em (Top)}}
\newcommand{\figbottom}{{\em (Bottom)}}
\newcommand{\captiona}{{\em (a)}}
\newcommand{\captionb}{{\em (b)}}
\newcommand{\captionc}{{\em (c)}}
\newcommand{\captiond}{{\em (d)}}

% Highlight a newly defined term
\newcommand{\newterm}[1]{{\bf #1}}


% Figure reference, lower-case.
\def\figref#1{figure~\ref{#1}}
% Figure reference, capital. For start of sentence
\def\Figref#1{Figure~\ref{#1}}
\def\twofigref#1#2{figures \ref{#1} and \ref{#2}}
\def\quadfigref#1#2#3#4{figures \ref{#1}, \ref{#2}, \ref{#3} and \ref{#4}}
% Section reference, lower-case.
\def\secref#1{section~\ref{#1}}
% Section reference, capital.
\def\Secref#1{Section~\ref{#1}}
% Reference to two sections.
\def\twosecrefs#1#2{sections \ref{#1} and \ref{#2}}
% Reference to three sections.
\def\secrefs#1#2#3{sections \ref{#1}, \ref{#2} and \ref{#3}}
% Reference to an equation, lower-case.
\def\eqref#1{equation~\ref{#1}}
% Reference to an equation, upper case
\def\Eqref#1{Equation~\ref{#1}}
% A raw reference to an equation---avoid using if possible
\def\plaineqref#1{\ref{#1}}
% Reference to a chapter, lower-case.
\def\chapref#1{chapter~\ref{#1}}
% Reference to an equation, upper case.
\def\Chapref#1{Chapter~\ref{#1}}
% Reference to a range of chapters
\def\rangechapref#1#2{chapters\ref{#1}--\ref{#2}}
% Reference to an algorithm, lower-case.
\def\algref#1{algorithm~\ref{#1}}
% Reference to an algorithm, upper case.
\def\Algref#1{Algorithm~\ref{#1}}
\def\twoalgref#1#2{algorithms \ref{#1} and \ref{#2}}
\def\Twoalgref#1#2{Algorithms \ref{#1} and \ref{#2}}
% Reference to a part, lower case
\def\partref#1{part~\ref{#1}}
% Reference to a part, upper case
\def\Partref#1{Part~\ref{#1}}
\def\twopartref#1#2{parts \ref{#1} and \ref{#2}}

\def\ceil#1{\lceil #1 \rceil}
\def\floor#1{\lfloor #1 \rfloor}
\def\1{\bm{1}}
\newcommand{\train}{\mathcal{D}}
\newcommand{\valid}{\mathcal{D_{\mathrm{valid}}}}
\newcommand{\test}{\mathcal{D_{\mathrm{test}}}}

\def\eps{{\epsilon}}


% Random variables
\def\reta{{\textnormal{$\eta$}}}
\def\ra{{\textnormal{a}}}
\def\rb{{\textnormal{b}}}
\def\rc{{\textnormal{c}}}
\def\rd{{\textnormal{d}}}
\def\re{{\textnormal{e}}}
\def\rf{{\textnormal{f}}}
\def\rg{{\textnormal{g}}}
\def\rh{{\textnormal{h}}}
\def\ri{{\textnormal{i}}}
\def\rj{{\textnormal{j}}}
\def\rk{{\textnormal{k}}}
\def\rl{{\textnormal{l}}}
% rm is already a command, just don't name any random variables m
\def\rn{{\textnormal{n}}}
\def\ro{{\textnormal{o}}}
\def\rp{{\textnormal{p}}}
\def\rq{{\textnormal{q}}}
\def\rr{{\textnormal{r}}}
\def\rs{{\textnormal{s}}}
\def\rt{{\textnormal{t}}}
\def\ru{{\textnormal{u}}}
\def\rv{{\textnormal{v}}}
\def\rw{{\textnormal{w}}}
\def\rx{{\textnormal{x}}}
\def\ry{{\textnormal{y}}}
\def\rz{{\textnormal{z}}}

% Random vectors
\def\rvepsilon{{\mathbf{\epsilon}}}
\def\rvtheta{{\mathbf{\theta}}}
\def\rva{{\mathbf{a}}}
\def\rvb{{\mathbf{b}}}
\def\rvc{{\mathbf{c}}}
\def\rvd{{\mathbf{d}}}
\def\rve{{\mathbf{e}}}
\def\rvf{{\mathbf{f}}}
\def\rvg{{\mathbf{g}}}
\def\rvh{{\mathbf{h}}}
\def\rvu{{\mathbf{i}}}
\def\rvj{{\mathbf{j}}}
\def\rvk{{\mathbf{k}}}
\def\rvl{{\mathbf{l}}}
\def\rvm{{\mathbf{m}}}
\def\rvn{{\mathbf{n}}}
\def\rvo{{\mathbf{o}}}
\def\rvp{{\mathbf{p}}}
\def\rvq{{\mathbf{q}}}
\def\rvr{{\mathbf{r}}}
\def\rvs{{\mathbf{s}}}
\def\rvt{{\mathbf{t}}}
\def\rvu{{\mathbf{u}}}
\def\rvv{{\mathbf{v}}}
\def\rvw{{\mathbf{w}}}
\def\rvx{{\mathbf{x}}}
\def\rvy{{\mathbf{y}}}
\def\rvz{{\mathbf{z}}}

% Elements of random vectors
\def\erva{{\textnormal{a}}}
\def\ervb{{\textnormal{b}}}
\def\ervc{{\textnormal{c}}}
\def\ervd{{\textnormal{d}}}
\def\erve{{\textnormal{e}}}
\def\ervf{{\textnormal{f}}}
\def\ervg{{\textnormal{g}}}
\def\ervh{{\textnormal{h}}}
\def\ervi{{\textnormal{i}}}
\def\ervj{{\textnormal{j}}}
\def\ervk{{\textnormal{k}}}
\def\ervl{{\textnormal{l}}}
\def\ervm{{\textnormal{m}}}
\def\ervn{{\textnormal{n}}}
\def\ervo{{\textnormal{o}}}
\def\ervp{{\textnormal{p}}}
\def\ervq{{\textnormal{q}}}
\def\ervr{{\textnormal{r}}}
\def\ervs{{\textnormal{s}}}
\def\ervt{{\textnormal{t}}}
\def\ervu{{\textnormal{u}}}
\def\ervv{{\textnormal{v}}}
\def\ervw{{\textnormal{w}}}
\def\ervx{{\textnormal{x}}}
\def\ervy{{\textnormal{y}}}
\def\ervz{{\textnormal{z}}}

% Random matrices
\def\rmA{{\mathbf{A}}}
\def\rmB{{\mathbf{B}}}
\def\rmC{{\mathbf{C}}}
\def\rmD{{\mathbf{D}}}
\def\rmE{{\mathbf{E}}}
\def\rmF{{\mathbf{F}}}
\def\rmG{{\mathbf{G}}}
\def\rmH{{\mathbf{H}}}
\def\rmI{{\mathbf{I}}}
\def\rmJ{{\mathbf{J}}}
\def\rmK{{\mathbf{K}}}
\def\rmL{{\mathbf{L}}}
\def\rmM{{\mathbf{M}}}
\def\rmN{{\mathbf{N}}}
\def\rmO{{\mathbf{O}}}
\def\rmP{{\mathbf{P}}}
\def\rmQ{{\mathbf{Q}}}
\def\rmR{{\mathbf{R}}}
\def\rmS{{\mathbf{S}}}
\def\rmT{{\mathbf{T}}}
\def\rmU{{\mathbf{U}}}
\def\rmV{{\mathbf{V}}}
\def\rmW{{\mathbf{W}}}
\def\rmX{{\mathbf{X}}}
\def\rmY{{\mathbf{Y}}}
\def\rmZ{{\mathbf{Z}}}

% Elements of random matrices
\def\ermA{{\textnormal{A}}}
\def\ermB{{\textnormal{B}}}
\def\ermC{{\textnormal{C}}}
\def\ermD{{\textnormal{D}}}
\def\ermE{{\textnormal{E}}}
\def\ermF{{\textnormal{F}}}
\def\ermG{{\textnormal{G}}}
\def\ermH{{\textnormal{H}}}
\def\ermI{{\textnormal{I}}}
\def\ermJ{{\textnormal{J}}}
\def\ermK{{\textnormal{K}}}
\def\ermL{{\textnormal{L}}}
\def\ermM{{\textnormal{M}}}
\def\ermN{{\textnormal{N}}}
\def\ermO{{\textnormal{O}}}
\def\ermP{{\textnormal{P}}}
\def\ermQ{{\textnormal{Q}}}
\def\ermR{{\textnormal{R}}}
\def\ermS{{\textnormal{S}}}
\def\ermT{{\textnormal{T}}}
\def\ermU{{\textnormal{U}}}
\def\ermV{{\textnormal{V}}}
\def\ermW{{\textnormal{W}}}
\def\ermX{{\textnormal{X}}}
\def\ermY{{\textnormal{Y}}}
\def\ermZ{{\textnormal{Z}}}

% Vectors
\def\vzero{{\bm{0}}}
\def\vone{{\bm{1}}}
\def\vmu{{\bm{\mu}}}
\def\vtheta{{\bm{\theta}}}
\def\va{{\bm{a}}}
\def\vb{{\bm{b}}}
\def\vc{{\bm{c}}}
\def\vd{{\bm{d}}}
\def\ve{{\bm{e}}}
\def\vf{{\bm{f}}}
\def\vg{{\bm{g}}}
\def\vh{{\bm{h}}}
\def\vi{{\bm{i}}}
\def\vj{{\bm{j}}}
\def\vk{{\bm{k}}}
\def\vl{{\bm{l}}}
\def\vm{{\bm{m}}}
\def\vn{{\bm{n}}}
\def\vo{{\bm{o}}}
\def\vp{{\bm{p}}}
\def\vq{{\bm{q}}}
\def\vr{{\bm{r}}}
\def\vs{{\bm{s}}}
\def\vt{{\bm{t}}}
\def\vu{{\bm{u}}}
\def\vv{{\bm{v}}}
\def\vw{{\bm{w}}}
\def\vx{{\bm{x}}}
\def\vy{{\bm{y}}}
\def\vz{{\bm{z}}}

% Elements of vectors
\def\evalpha{{\alpha}}
\def\evbeta{{\beta}}
\def\evepsilon{{\epsilon}}
\def\evlambda{{\lambda}}
\def\evomega{{\omega}}
\def\evmu{{\mu}}
\def\evpsi{{\psi}}
\def\evsigma{{\sigma}}
\def\evtheta{{\theta}}
\def\eva{{a}}
\def\evb{{b}}
\def\evc{{c}}
\def\evd{{d}}
\def\eve{{e}}
\def\evf{{f}}
\def\evg{{g}}
\def\evh{{h}}
\def\evi{{i}}
\def\evj{{j}}
\def\evk{{k}}
\def\evl{{l}}
\def\evm{{m}}
\def\evn{{n}}
\def\evo{{o}}
\def\evp{{p}}
\def\evq{{q}}
\def\evr{{r}}
\def\evs{{s}}
\def\evt{{t}}
\def\evu{{u}}
\def\evv{{v}}
\def\evw{{w}}
\def\evx{{x}}
\def\evy{{y}}
\def\evz{{z}}

% Matrix
\def\mA{{\bm{A}}}
\def\mB{{\bm{B}}}
\def\mC{{\bm{C}}}
\def\mD{{\bm{D}}}
\def\mE{{\bm{E}}}
\def\mF{{\bm{F}}}
\def\mG{{\bm{G}}}
\def\mH{{\bm{H}}}
\def\mI{{\bm{I}}}
\def\mJ{{\bm{J}}}
\def\mK{{\bm{K}}}
\def\mL{{\bm{L}}}
\def\mM{{\bm{M}}}
\def\mN{{\bm{N}}}
\def\mO{{\bm{O}}}
\def\mP{{\bm{P}}}
\def\mQ{{\bm{Q}}}
\def\mR{{\bm{R}}}
\def\mS{{\bm{S}}}
\def\mT{{\bm{T}}}
\def\mU{{\bm{U}}}
\def\mV{{\bm{V}}}
\def\mW{{\bm{W}}}
\def\mX{{\bm{X}}}
\def\mY{{\bm{Y}}}
\def\mZ{{\bm{Z}}}
\def\mBeta{{\bm{\beta}}}
\def\mPhi{{\bm{\Phi}}}
\def\mLambda{{\bm{\Lambda}}}
\def\mSigma{{\bm{\Sigma}}}

% Tensor
\DeclareMathAlphabet{\mathsfit}{\encodingdefault}{\sfdefault}{m}{sl}
\SetMathAlphabet{\mathsfit}{bold}{\encodingdefault}{\sfdefault}{bx}{n}
\newcommand{\tens}[1]{\bm{\mathsfit{#1}}}
\def\tA{{\tens{A}}}
\def\tB{{\tens{B}}}
\def\tC{{\tens{C}}}
\def\tD{{\tens{D}}}
\def\tE{{\tens{E}}}
\def\tF{{\tens{F}}}
\def\tG{{\tens{G}}}
\def\tH{{\tens{H}}}
\def\tI{{\tens{I}}}
\def\tJ{{\tens{J}}}
\def\tK{{\tens{K}}}
\def\tL{{\tens{L}}}
\def\tM{{\tens{M}}}
\def\tN{{\tens{N}}}
\def\tO{{\tens{O}}}
\def\tP{{\tens{P}}}
\def\tQ{{\tens{Q}}}
\def\tR{{\tens{R}}}
\def\tS{{\tens{S}}}
\def\tT{{\tens{T}}}
\def\tU{{\tens{U}}}
\def\tV{{\tens{V}}}
\def\tW{{\tens{W}}}
\def\tX{{\tens{X}}}
\def\tY{{\tens{Y}}}
\def\tZ{{\tens{Z}}}


% Graph
\def\gA{{\mathcal{A}}}
\def\gB{{\mathcal{B}}}
\def\gC{{\mathcal{C}}}
\def\gD{{\mathcal{D}}}
\def\gE{{\mathcal{E}}}
\def\gF{{\mathcal{F}}}
\def\gG{{\mathcal{G}}}
\def\gH{{\mathcal{H}}}
\def\gI{{\mathcal{I}}}
\def\gJ{{\mathcal{J}}}
\def\gK{{\mathcal{K}}}
\def\gL{{\mathcal{L}}}
\def\gM{{\mathcal{M}}}
\def\gN{{\mathcal{N}}}
\def\gO{{\mathcal{O}}}
\def\gP{{\mathcal{P}}}
\def\gQ{{\mathcal{Q}}}
\def\gR{{\mathcal{R}}}
\def\gS{{\mathcal{S}}}
\def\gT{{\mathcal{T}}}
\def\gU{{\mathcal{U}}}
\def\gV{{\mathcal{V}}}
\def\gW{{\mathcal{W}}}
\def\gX{{\mathcal{X}}}
\def\gY{{\mathcal{Y}}}
\def\gZ{{\mathcal{Z}}}

% Sets
\def\sA{{\mathbb{A}}}
\def\sB{{\mathbb{B}}}
\def\sC{{\mathbb{C}}}
\def\sD{{\mathbb{D}}}
% Don't use a set called E, because this would be the same as our symbol
% for expectation.
\def\sF{{\mathbb{F}}}
\def\sG{{\mathbb{G}}}
\def\sH{{\mathbb{H}}}
\def\sI{{\mathbb{I}}}
\def\sJ{{\mathbb{J}}}
\def\sK{{\mathbb{K}}}
\def\sL{{\mathbb{L}}}
\def\sM{{\mathbb{M}}}
\def\sN{{\mathbb{N}}}
\def\sO{{\mathbb{O}}}
\def\sP{{\mathbb{P}}}
\def\sQ{{\mathbb{Q}}}
\def\sR{{\mathbb{R}}}
\def\sS{{\mathbb{S}}}
\def\sT{{\mathbb{T}}}
\def\sU{{\mathbb{U}}}
\def\sV{{\mathbb{V}}}
\def\sW{{\mathbb{W}}}
\def\sX{{\mathbb{X}}}
\def\sY{{\mathbb{Y}}}
\def\sZ{{\mathbb{Z}}}

% Entries of a matrix
\def\emLambda{{\Lambda}}
\def\emA{{A}}
\def\emB{{B}}
\def\emC{{C}}
\def\emD{{D}}
\def\emE{{E}}
\def\emF{{F}}
\def\emG{{G}}
\def\emH{{H}}
\def\emI{{I}}
\def\emJ{{J}}
\def\emK{{K}}
\def\emL{{L}}
\def\emM{{M}}
\def\emN{{N}}
\def\emO{{O}}
\def\emP{{P}}
\def\emQ{{Q}}
\def\emR{{R}}
\def\emS{{S}}
\def\emT{{T}}
\def\emU{{U}}
\def\emV{{V}}
\def\emW{{W}}
\def\emX{{X}}
\def\emY{{Y}}
\def\emZ{{Z}}
\def\emSigma{{\Sigma}}

% entries of a tensor
% Same font as tensor, without \bm wrapper
\newcommand{\etens}[1]{\mathsfit{#1}}
\def\etLambda{{\etens{\Lambda}}}
\def\etA{{\etens{A}}}
\def\etB{{\etens{B}}}
\def\etC{{\etens{C}}}
\def\etD{{\etens{D}}}
\def\etE{{\etens{E}}}
\def\etF{{\etens{F}}}
\def\etG{{\etens{G}}}
\def\etH{{\etens{H}}}
\def\etI{{\etens{I}}}
\def\etJ{{\etens{J}}}
\def\etK{{\etens{K}}}
\def\etL{{\etens{L}}}
\def\etM{{\etens{M}}}
\def\etN{{\etens{N}}}
\def\etO{{\etens{O}}}
\def\etP{{\etens{P}}}
\def\etQ{{\etens{Q}}}
\def\etR{{\etens{R}}}
\def\etS{{\etens{S}}}
\def\etT{{\etens{T}}}
\def\etU{{\etens{U}}}
\def\etV{{\etens{V}}}
\def\etW{{\etens{W}}}
\def\etX{{\etens{X}}}
\def\etY{{\etens{Y}}}
\def\etZ{{\etens{Z}}}

% The true underlying data generating distribution
\newcommand{\pdata}{p_{\rm{data}}}
% The empirical distribution defined by the training set
\newcommand{\ptrain}{\hat{p}_{\rm{data}}}
\newcommand{\Ptrain}{\hat{P}_{\rm{data}}}
% The model distribution
\newcommand{\pmodel}{p_{\rm{model}}}
\newcommand{\Pmodel}{P_{\rm{model}}}
\newcommand{\ptildemodel}{\tilde{p}_{\rm{model}}}
% Stochastic autoencoder distributions
\newcommand{\pencode}{p_{\rm{encoder}}}
\newcommand{\pdecode}{p_{\rm{decoder}}}
\newcommand{\precons}{p_{\rm{reconstruct}}}

\newcommand{\laplace}{\mathrm{Laplace}} % Laplace distribution

\newcommand{\E}{\mathbb{E}}
\newcommand{\Ls}{\mathcal{L}}
\newcommand{\R}{\mathbb{R}}
\newcommand{\emp}{\tilde{p}}
\newcommand{\lr}{\alpha}
\newcommand{\reg}{\lambda}
\newcommand{\rect}{\mathrm{rectifier}}
\newcommand{\softmax}{\mathrm{softmax}}
\newcommand{\sigmoid}{\sigma}
\newcommand{\softplus}{\zeta}
\newcommand{\KL}{D_{\mathrm{KL}}}
\newcommand{\Var}{\mathrm{Var}}
\newcommand{\standarderror}{\mathrm{SE}}
\newcommand{\Cov}{\mathrm{Cov}}
% Wolfram Mathworld says $L^2$ is for function spaces and $\ell^2$ is for vectors
% But then they seem to use $L^2$ for vectors throughout the site, and so does
% wikipedia.
\newcommand{\normlzero}{L^0}
\newcommand{\normlone}{L^1}
\newcommand{\normltwo}{L^2}
\newcommand{\normlp}{L^p}
\newcommand{\normmax}{L^\infty}

\newcommand{\parents}{Pa} % See usage in notation.tex. Chosen to match Daphne's book.

\DeclareMathOperator*{\argmax}{arg\,max}
\DeclareMathOperator*{\argmin}{arg\,min}

\DeclareMathOperator{\sign}{sign}
\DeclareMathOperator{\Tr}{Tr}
\let\ab\allowbreak

\usepackage{hyperref}
\usepackage{url}
\usepackage{authblk}
\usepackage{booktabs}       % professional-quality tables
\usepackage{amsfonts}       % blackboard math symbols
\usepackage{nicefrac}       % compact symbols for 1/2, etc.
\usepackage{microtype}      % microtypography
\usepackage{xcolor}         % colors
\usepackage{amsmath,amssymb}
\usepackage{amsthm}
\usepackage {cancel}
\usepackage{autobreak}
\usepackage{algorithm}
\usepackage[noend]{algpseudocode}
\usepackage{algorithmicx}
\usepackage{bm}
\usepackage{paralist}
\usepackage{here}
\usepackage{multirow}
\usepackage{wrapfig}
\usepackage{subcaption}
\usepackage{caption}
\captionsetup[subfigure]{justification=centering}
\usepackage{fancyhdr}       % header
\usepackage{graphicx}       % graphics
\usepackage{tikz}
\usepackage{mathrsfs}
\usepackage{mathtools}
\usetikzlibrary{patterns}
\usepackage{changepage}
\usepackage[capitalise]{cleveref}

\makeatletter
\usepackage{xspace} 
\def\@onedot{\ifx\@let@token.\else.\null\fi\xspace}
\DeclareRobustCommand\onedot{\futurelet\@let@token\@onedot}
\DeclarePairedDelimiterX\setc[2]{\{}{\}}{\,#1 \;\delimsize\vert\; #2\,}

\def\eg{\emph{e.g}\onedot}
\def\Eg{\emph{E.g}\onedot}
\def\ie{\emph{i.e}\onedot}
\def\Ie{\emph{I.e}\onedot}
\def\cf{\emph{cf}\onedot}
\def\Cf{\emph{Cf}\onedot}
\def\etc{\emph{etc}\onedot}
\def\vs{\emph{vs}\onedot}
\def\wrt{w.r.t\onedot}
\def\dof{d.o.f\onedot}
\def\aka{a.k.a\onedot}
\def\iid{i.i.d\onedot}
\def\etal{\emph{et al}\onedot}

\title{Understanding Generalization in Physics Informed Models through Affine Variety Dimensions}

\author[1]{Takeshi Koshizuka}
\author[1]{Issei Sato}
\affil[1]{Department of Computer Science, The University of Tokyo}
\affil[1]{\{koshizuka-takeshi938444, sato\}@g.ecc.u-tokyo.ac.jp}

\newcommand{\fix}{\marginpar{FIX}}
\newcommand{\new}{\marginpar{NEW}}

\renewcommand{\algorithmicrequire}{\textbf{Input:}}
\renewcommand{\algorithmicensure}{\textbf{Output:}}

\usepackage{thmtools}
\usepackage{thm-restate}

\theoremstyle{plain}
\newtheorem{theorem}{Theorem}[section]
\newtheorem{proposition}[theorem]{Proposition}
\newtheorem{lemma}[theorem]{Lemma}
\newtheorem{corollary}[theorem]{Corollary}
\theoremstyle{definition}
\newtheorem{definition}[theorem]{Definition}
\newtheorem{assumption}[theorem]{Assumption}
\theoremstyle{remark}
\newtheorem{remark}[theorem]{Remark}

\begin{document}
\maketitle
\begin{abstract}
Physics-informed machine learning is gaining significant traction for enhancing statistical performance and sample efficiency through the integration of physical knowledge. However, current theoretical analyses often presume complete prior knowledge in non-hybrid settings, overlooking the crucial integration of observational data, and are frequently limited to linear systems, unlike the prevalent nonlinear nature of many real-world applications.
To address these limitations, we introduce a unified residual form that unifies collocation and variational methods, enabling the incorporation of incomplete and complex physical constraints in hybrid learning settings.
Within this formulation, we establish that the generalization performance of physics-informed regression in such hybrid settings is governed by the dimension of the affine variety associated with the physical constraint, rather than by the number of parameters. This enables a unified analysis that is applicable to both linear and nonlinear equations. We also present a method to approximate this dimension and provide experimental validation of our theoretical findings.
\end{abstract}
\section{Introduction}
In recent years, physics-informed machine learning (PIML) has garnered significant attention~\citep{rai2020driven,karniadakis2021physics,cuomo2022scientific,hao2022physics}. PIML represents a hybrid approach that integrates physical knowledge into machine learning models for tasks involving physical phenomena. These hybrid models can leverage inherent physical structures, such as differential equations~\citep{raissi2019physics}, conservation laws~\citep{jagtap2020conservative}, and symmetries~\citep{akhound2024lie}, as inductive biases. This integration has the potential to enhance both sample efficiency and generalization capabilities. These models have been empirically applied to a wide range of phenomena, with successful applications observed in areas including thrombus material properties \citep{yin2021non}, fluid dynamics \citep{cai2021flow,jin2021nsfnets}, turbulence \citep{wang2020towards}, and heat transfer problems~\citep{cai2021physics}.

Despite these empirical successes, the theoretical analysis of PIML remains underdeveloped, which potentially undermines the reliability of these methods. Notably, in practical scenarios, prior knowledge of the governing differential equations, particularly their source terms or boundary conditions, is often incomplete. Consequently, learning frequently involves a hybrid approach of fitting to actual observational data alongside incorporating physical constraints. However, much of the existing theoretical research focuses on settings where complete prior knowledge of the differential equations is assumed. Furthermore, the scope of existing analyses is often limited to linear differential equations or systems exhibiting strong regularity, creating a gap between theory and application.

To bridge this gap, we propose a versatile analytical framework for physics-informed regression using linear hypothesis classes over nonlinear features in hybrid settings. Our key idea is to formulate the differential equation constraints by introducing a Unified Residual Form. This form, defined on a finite set of trial functions and a measure, provides a practical approximation of physical constraints. This formulation unifies the collocation-based constraints, used in Physics-Informed Neural Networks (PINNs)~\citep{raissi2019physics}, and the standard unified residual form constraints, used in the variational and finite element methods. Within this framework, the learning weights of a physics-informed linear regressor are shown to be defined on an affine variety associated with the differential equation. Crucially, our analysis aims to elucidate the impact of incorporating physical prior knowledge on the generalization capacity of these models. We then establish that the generalization capacity of these models is determined by the dimension of this affine variety, rather than by the number of parameters. This novel perspective enables a unified analysis applicable to various equations, including nonlinear ones. To support our theoretical findings, we introduce a method for approximately calculating the dimension of the affine variety and provide extensive experimental validation. Our results illustrate that even in scenarios with a large number of parameters relative to the amount of data, the incorporation of physical structure reduces the intrinsic dimension of the hypothesis space, thereby mitigating overfitting and corroborating our theoretical claims.

% The remainder of this paper is organized as follows. In~\cref{sec:analysis}, we outline the problem setup and present our main theoretical results, including a minimax risk analysis that underscores the role of the dimension of the affine variety. In~\cref{sec:dim-variety}, we discuss the dimension of affine variety especially in the context of nonlinear operators and introduce methods for their approximate calculation. \Cref{sec:exp} provides experimental evidence supporting our theoretical claims, demonstrating the practical advantages of incorporating physical structures in machine learning models. 

\section{Related Work}
Since the seminal work by~\citet{raissi2019physics} on PINNs, PIML has rapidly emerged as a significant field of study. This area has been comprehensively surveyed in the literature by~\citep{rai2020driven,karniadakis2021physics,cuomo2022scientific,hao2022physics}. Leveraging the high function approximation capabilities of neural networks~\citep{hornik1989multilayer,kutyniok2022theoretical,de2022error}, these models have been employed as versatile surrogates for solving various equations. In contrast, linear models are also used because of their interpretability, consistency with classical numerical solvers~\citep{arnone2022some,ferraccioli2022some}, and the close relationship between Partial Differential Equations (PDEs) and kernel methods~\citep{schaback2006kernel,chen2021solving,long2022autoip,dalton2024physics,doumeche2024physics2}. Recently, methods that exploit underlying conservation laws~\citep{jagtap2020conservative,hu2022extended} and symmetries~\citep{akhound2024lie,dalton2024physics}, in addition to the equations themselves, have also been developed.

Recent studies have made advances in the theoretical understanding of PINNs. \citet{shin2020convergence} rigorously showed that the minimizer of the PINN loss converges to the strong solution as the data size approaches infinity for linear elliptic and parabolic PDEs under certain conditions. These findings were extended by~\citet{shin2023error} into a general framework applicable to broader linear problems, with the loss function formulated in both strong and variational forms. \citet{mishra2022estimates,mishra2023estimates} use the stability properties of the underlying PDEs to derive upper bounds on the generalization error of PINNs. Subsequent research has applied this analytical framework to various specific equations~\citep{bai2021physics,mishra2021physics}. However, studies explicitly addressing the impact of physical structure on generalization capabilities are still limited. \citet{arnone2022some} proved that for second-order elliptic PDEs, the physics-informed linear estimator using a finite element basis converges at a rate surpassing the Sobolev minimax rate. \citet{doumeche2024physics} quantified the generalization capacity of the physics-informed estimator for general linear PDEs using the concept of effective dimension~\citep{caponnetto2007optimal}, a well-known metric in the analysis of the kernel method. The effects of incorporating the structures of nonlinear complex equations, as well as conservation laws and symmetries, into models on generalization, have yet to be thoroughly analyzed.

\section{Minimax risk Analysis}
\label{sec:analysis}
In this section, we explain how introducing physical structures can improve the generalization capacity of linear models. In~\cref{sec:problem-setup}, we outline the problem setup. In~\cref{sec:main-theorem}, we perform a minimax risk analysis, showing that the generalization capacity is mainly determined by the dimension of the affine variety. In~\cref{sec:linop}, we show that our theory aligns with existing theories on linear operators. Notations are summarized in~\cref{app:notation}. 

\subsection{Problem Setup}
\label{sec:problem-setup}
We consider the regression problem, which aims to learn the unknown function \( f^* \colon \mathbb{R}^m \to \mathbb{R} \) that satisfies the differential equation. We have a dataset consisting of \( n \) observations, denoted as \( \{(x_i, y_i)\}_{i=1}^n \), where \( x_i \in \Omega \cup \partial \Omega \) represents the input within the domain \( \Omega \subseteq \mathbb{R}^m \) or the boundary $\partial \Omega$ and \( y_i \in \mathbb{R} \) represents the corresponding output. Observations are sampled independently from a probability distribution \(\mathcal{P}\) on the domain \(\Omega \cup \partial \Omega \times \mathbb{R}\). The relationship between the observations and the true function can be expressed as:
\begin{align*}
    y_i = f^*(x_i) + \varepsilon_i,\ \varepsilon_i \sim \mathcal{N}(0, \sigma^2),
\end{align*}
where \( \varepsilon_i \) represents normally distributed noise with mean zero and variance \( \sigma^2 \). 
The target function \(f^* \) is the solution of the differential equation, \ie, \( \mathscr{D}[f^*] = 0 \) for a given operator \( \mathscr{D} \colon L^2(\Omega) \to L^2(\Omega) \), where $L^2(\Omega)$ denotes the space of square-integrable functions on a domain $\Omega\subseteq \mathbb{R}^m$. For a more detailed background on the problem setting, please refer to \cref{app:detail_background}. 
%The values of the boundary conditions are unknown and are learned from fitting the observational data.

\paragraph{Unified Residual Form:} To formulate incomplete prior knowledge of the governing differential equations, we first introduce a unified residual form, which captures physical constraints in an integrated or averaged sense. Such formulations naturally arise in variational and finite element methods, and are particularly well-suited to hybrid settings where only partial physical supervision is available. 
Formally, let $\mathcal{T} \;\coloneqq\; \left\{(\psi_k,\mu_k)\right\}_{k=1}^K$ be a finite collection of trial functions and measure pairs, where each \(\psi_k:\mathbb{R}^m\to\mathbb{R}\) is a smooth trial function and \(\mu_k: \Sigma \to \mathbb{R}\) is a measure on the \(\sigma\)‐algebra \(\Sigma\) over the domain \(\Omega\). Then, the unified residual form of the knwon differential equation $\mathscr{D}$ is defined by
\begin{align*}
\langle \mathscr{D}[f],\psi_k\rangle_{\mu_k}
\;\coloneqq\;\int_{\Omega}\!\mathscr{D}[f](x)\,\psi_k(x)\,\mathrm{d}\mu_k(x)
\;=\;0,
\quad k=1,\dots,K.
\end{align*}

We then impose the differential equation constraint through the unified residual form defined above. The resulting physics‐informed regression problem reads:
\begin{gather}
  \label{eq:problem}
  \begin{aligned}
    \hat f_n
    &=\;\argmin_{f\in\mathcal{F}(\mathscr{D},\mathcal{T})}
      \frac{1}{n}\sum_{i=1}^n \left|y_i - f(x_i)\right|^2
      + \lambda_n \|f\|^2,\\
    \mathcal{F}(\mathscr{D},\mathcal{T})
    &\coloneqq
      \left\{\,f : \langle \mathscr{D}[f],\psi_k\rangle_{\mu_k}=0,\ 
      \forall(\psi_k,\mu_k)\in\mathcal{T}\right\},
  \end{aligned}
\end{gather}
where \(\lambda_n\) is a regularization parameter and \(\|\cdot\|\) denotes the standard \(L^2\) norm. This formulation relaxes the classical smoothness requirements while still leveraging physics-informed constraints via a unified measure-based approach: choosing Borel measures leads to an approximation of standard weak solutions, whereas choosing Dirac measures leads to an approximation of the strong-form residuals used in the PINN framework.

\paragraph{Physics-Informed Linear Regression (PILR) Setup:} Let $\mathcal{B}=\{\phi_j:\mathbb{R}^m\to\mathbb{R}\}_{j=1}^d$ be a fixed basis. Define the basis vector $\boldsymbol{\phi}(x) = [\phi_1(x), \phi_2(x), \dots, \phi_d(x)]^\top \in \mathbb{R}^d$, the design matrix $\mathbf{\Phi} = [\boldsymbol{\phi}(x_1), \boldsymbol{\phi}(x_2), \dots, \boldsymbol{\phi}(x_n)]^\top \in \mathbb{R}^{n \times d}$, and the target vector $\boldsymbol{y}=[y_1,\dots,y_n]^\top\in\mathbb{R}^n$.

The physics‐informed feasible set is
\begin{align}
\label{eq:def-V}
\mathcal{V}(\mathscr{D}, \mathcal{B}, \mathcal{T}) \coloneqq \left\{ \bm{w} \in \mathbb{R}^d: \langle \mathscr{D}\left[\bm{w}^{\top} \bm{\phi} \right], \psi_k \rangle_{\mu_k} = 0,\ \forall (\psi_k, \mu_k) \in \mathcal{T}, \phi_j \in \mathcal{B} \right\}.
\end{align}
The problem \cref{eq:problem} reduces to the physics-informed linear regression given by
\begin{gather}
\label{eq:lin-problem-ball}
\widehat{\boldsymbol{w}} = \underset{\boldsymbol{w} \in \mathcal{V}_R}{\arg \min } \frac{1}{n} \| \boldsymbol{y} - \mathbf{\Phi} \boldsymbol{w} \|_2^2,
\end{gather}
where $\mathcal{V}_R = \mathcal{V}(\mathscr{D}, \mathcal{B}, \mathcal{T}) \cap \mathbb{B}_2(R)$ is the affine variety constrained by the $\ell_2$-ball $\mathbb{B}_2(R)$ with radius $R > 0$ and \(\|\cdot\|_2\) is the \(\ell_2\)-norm.  

The set of coefficients \( \mathcal{V} \) constitutes \textit{an affine variety} as it represents the set of solutions to the \(K\) polynomial equations in the \( d \) variables with real coefficients. For example, when \( m = 1 \) and \( \mathscr{D}[f] = f \cdot \frac{\mathrm{d}}{\mathrm{d}x} f \), the affine variety \( \mathcal{V} \) is defined by the solution set of the polynomial equations \( p_k(\bm{w}) = \sum_{j, j'=1}^d \langle \left(\frac{\mathrm{d}}{\mathrm{d}x} \phi_j \right)\phi_{j'}, \psi_k \rangle_{\mu_k} w_j w_{j'} = 0 \) for \( k = 1, \ldots, K \). 
We perform minimax risk analysis based on the dimension \( d_{\mathcal{V}} \) of this affine variety because the affine variety \( \mathcal{V} \) is crucial to determine the size of the intrinsic hypothesis space.

\textbf{Minimax risk:} The goal of our analysis is to obtain the upper bound of the minimax risk for PILR in \cref{eq:lin-problem-ball}, which is defined by
\begin{gather}
    \label{eq:def-minmax-risk}
    \min_{\hat{\bm{w}}} \max_{\bm{w}^* \in \mathcal{V}_R} 
    \|\hat{\bm{w}} - \bm{w}^*\|_2^2,
\end{gather}
Here, $\bm{w}^* \in \mathcal{V}_R$ represents the optimal weight vector. The corresponding optimal hypothesis $f_{\bm{w}^*} = \bm{w}^{*\top} \bm{\phi}$ within our hypothesis space $\mathcal{H} = \{ \bm{w}^\top \bm{\phi} : \boldsymbol{w} \in \mathcal{V}_R \}$ is defined as the best approximation of the true function $f^*$: $f_{\bm{w}^*} = \bm{w}^{*\top} \bm{\phi} = \argmin_{f_{\bm{w}} \in \mathcal{H}} \|f_{\bm{w}} - f^*\|^2$.

\textbf{We strongly recommend referring to the example in~\cref{sec:strong-sol} to intuitively understand our problem setting.}

\subsection{Main Theorem}
\label{sec:main-theorem}
In this section, we present an upper bound on the minimax risk for PILR. 
The bound is interpretable and sufficiently sharp, revealing how physical constraints reduce hypothesis complexity and enhance generalization.
We begin by stating the definition and assumptions underpinning our analysis.
\begin{definition}[$(\beta,d_{\mathcal V})$-regular set]
An affine variety \(V \subseteq \mathbb{R}^d\) is called a \((\beta, d_{\mathcal V})\)-regular set if the following conditions hold: (1) For almost all affine subspaces \(L \subseteq \mathbb{R}^d\) of dimension \(d_L\) satisfying \(d - d_L \leq d_{\mathcal V}\), the intersection \(V \cap L\) has at most \(\beta\) path-connected components. (2) For almost all affine subspaces \(L \subseteq \mathbb{R}^d\) of dimension \(d_L\) with \(d - d_L > d_{\mathcal V}\), the intersection \(V \cap L\) is empty. See~\cref{app:regularity_of_affinev} for illustrative explanations.
\end{definition}
\begin{assumption}[Boundedness of basis functions]
\label{asm:bounded_phi}
For the basis function $\bm{\phi} = [\phi_1, \ldots, \phi_d]^{\top}$, where $\phi_j \in \mathcal{B}$, assume that there exists a positive constant $M$ such that $\| \bm{\phi}(x) \|_2 \leq M$ for all $x \in \Omega$.
\end{assumption}
\begin{assumption}
\label{asm:lowereig}
Assume there exists a constant $\eta > 0$ such that $ \frac{1}{\sqrt{n}} \|\bm{\Phi}\bm{w}\|_2 \geq \sqrt{\eta} \|\bm{w}\|_2 $ for all $\bm{w} \in \mathbb{B}_2(2R)$.
\end{assumption}
\begin{assumption}[Stability of estimator]
\label{asm:stability_est}
Assume there exists a constant $\Gamma > 1$ such that $\| \hat{\bm{w}}_1 - \hat{\bm{w}}_2 \|_2 \leq (\Gamma - 1) \| \bm{w}_1^* - \bm{w}_2^* \|_2$, for the estimators $\hat{\bm{w}}_1$ and $\hat{\bm{w}}_2$ of the optimal weights $\bm{w}_1^*$ and $\bm{w}_2^*$, respectively.
\end{assumption}

Next, we present the upper bound on the minimax risk. The complete proof is provided in~\cref{app:proof_thm:main}.
\begin{restatable}[Minimax Risk Bound]{theorem}{main}\label{thm:main}
Let $\mathcal{V}(\mathscr{D}, \mathcal{B}, \mathcal{T})$ be the $(\beta, d_{\mathcal{V}})$-regular affine variety defined in~\cref{eq:def-V}. Suppose Assumptions~\ref{asm:bounded_phi}-\ref{asm:stability_est} hold. Then, there exists a positive constant $C$, independent of $n$, $d_{\mathcal{V}}$, $d$, and $\beta$, such that for any $\delta \in (0, 1)$, with probability at least $1-\delta$, the minimax risk for PILR defined by~\cref{eq:def-minmax-risk} is bounded by
\begin{align}
    \label{eq:thm:main}
    \min_{\hat{\bm{w}}} \max_{\bm{w}^* \in \mathcal{V}_R} \|\hat{\bm{w}} - \bm{w}^*\|^2_2 \leq C \eta^{-1} \sigma M \Gamma R \left(\sqrt{\frac{d_{\mathcal{V}} \log(d_{\mathcal{V}} d)}{n}} + \sqrt{\frac{\log 2\beta}{n}} + 2\sqrt{\frac{\log(2/\delta)}{n}}\right).
\end{align}
\end{restatable}
\begin{proof}[Proof Sketch]
The proof proceeds in two steps. In the first step, we upper bound the minimax risk by the supremum of a sub-Gaussian random process defined over the metric space $(\mathcal{V}_R, \| \cdot \|_2)$. The second step utilizes Dudley's integral theorem, which bounds the supremum of the process by an integral involving its covering number, specifically: $\int_0^{\infty} \sqrt{\mathcal{N}(\mathcal{V}_R, \varepsilon, \|\cdot\|_2)} \, \mathrm{d}\varepsilon$. To apply Dudley's theorem effectively, we employ~\cref{lem:cov} to obtain an explicit upper bound for the covering number. Substituting this bound into Dudley's integral and performing the integration yields the desired high-probability minimax risk bound.
\end{proof}
% interpretation (d_V と \beta: ) 
\Cref{thm:main} demonstrates that the minimax risk is primarily governed by the intrinsic dimension \( d_{\mathcal{V}} \) of the affine variety \( \mathcal{V} \), rather than the ambient input dimension \( d \), particularly when the topological complexity parameter \( \beta \) is small. For comparison, standard least-squares estimation over an \(\ell_2\)-ball \( \mathbb{B}_2(R) \subset \mathbb{R}^d \) yields a minimax risk rate of order \( \mathcal{O}(\sqrt{d/n}) \), which is optimal for unconstrained linear regression in \( d \)-dimensional space. In contrast, our result shows that when \( d_{\mathcal{V}} \ll d \), incorporating physical structure into the hypothesis space through differential constraints significantly sharpens the risk rate, yielding improved generalization.

\paragraph{On the Role of \(\beta\).}
The parameter \( \beta \) captures the topological complexity of the affine variety and appears as a regularity constant in the generalization bound. Its upper bound can be estimated via the Petrovskii–Oleinik–Milnor inequality~\citep{Petrovskii1949,oleinik1951estimates,milnor1964betti}, which provides a bound on the sum of Betti numbers of a semialgebraic set. Specifically, if the variety \( \mathcal{V} \cap \mathbb{B}_2(R) \subset \mathbb{R}^d \) is defined by polynomial constraints \( \{ p_k(\bm{w}) \}_{k=1}^K \) of maximal degree \( \rho \), then it is \( (\rho(2\rho-1)^{d+1}, d_{\mathcal{V}}) \)-regular. This implies that as the degree \(\rho\) of the defining polynomials increases, the variety can exhibit more intricate topological features, such as additional holes and disconnected components.

\paragraph{How \(d_{\mathcal{V}}\) and \(\beta\) Arise from the Covering Argument.}
The minimax risk is bounded via Dudley’s entropy integral, which requires control over the covering number \( \mathcal{N}(\mathcal{V}_R, \varepsilon, \|\cdot\|_2) \). Following the geometric approach of~\citet{zhang2023covering}, the affine variety \( \mathcal{V} \subset \mathbb{R}^d \) is sliced using a family of linear subspaces \( \{L_s\}_{s \in \mathbb{N}} \), and each intersection \( \mathcal{V} \cap L_s \) is covered by Euclidean balls of radius \( \varepsilon \). The total covering is then given by
\begin{align*}
\mathcal{V} \subset \bigcup_s \bigcup_{v \in \mathcal{V} \cap L_s} \mathbb{B}_2(v;\varepsilon).
\end{align*}
In this construction, the intrinsic dimension \( d_{\mathcal{V}} \) controls the number of subspaces required to sufficiently cover \( \mathcal{V} \), while the parameter \( \beta \), corresponding to the sum of Betti numbers, governs the covering number of each individual section \( \mathcal{V} \cap L_s \). Topologically, \(\beta\) can be interpreted as quantifying the number of topological features (e.g., holes) in \( \mathcal{V} \), and thus reflects the local geometric complexity encountered within each subspace. For reference, the standard covering number of the Euclidean ball satisfies \( \mathcal{N}(\mathbb{B}_2(R), \varepsilon, \|\cdot\|_2) \leq (1+2R/\varepsilon)^d \), highlighting the advantage of replacing ambient-dimension dependence with complexity parameters intrinsic to the constraint set.

\paragraph{Key Insights.}
A central contribution of our analysis is its interpretability through the lens of intrinsic complexity measures. The dimension \( d_{\mathcal{V}} \) plays a role analogous to the VC dimension in classification~\citep{abu1989vapnik} or the pseudo-dimension in regression~\citep{pollard1990empirical}, serving as a proxy for the effective capacity of the hypothesis space. This dimensional viewpoint clarifies how the incorporation of physical constraints—via differential equation structure—can substantially reduce hypothesis complexity, even in high-dimensional ambient spaces. While this may come at the cost of slightly looser constants compared to minimax-optimal bounds, the resulting rate is still sharp enough to meaningfully capture the generalization benefit of physics-informed inductive bias. Empirical evidence supporting this theoretical advantage is presented in~\cref{sec:exp}, and an alternative analysis via Rademacher complexity is provided in~\cref{app:Rademacher}.

\paragraph{Effect of the Trial Function Set \(\mathcal{T}\).}
The set of trial functions \( \mathcal{T} \) encodes the imposed physical constraints, typically derived from a governing differential operator \( \mathscr{D} \). The cardinality \( K = |\mathcal{T}| \) quantifies the amount of physical knowledge embedded in the learning problem. Increasing the number of trial functions leads to a more restrictive constraint set, which geometrically corresponds to a lower-dimensional affine variety. Specifically, if \( \mathcal{T}_1 \subset \mathcal{T}_2 \), then it follows that
\begin{align*}
\mathcal{V}(\mathscr{D}, \mathcal{B}, \mathcal{T}_2) \subset \mathcal{V}(\mathscr{D}, \mathcal{B}, \mathcal{T}_1) \quad \Rightarrow \quad d_{\mathcal{V}(\mathscr{D}, \mathcal{B}, \mathcal{T}_2)} \leq d_{\mathcal{V}(\mathscr{D}, \mathcal{B}, \mathcal{T}_1)},
\end{align*}
which highlights how adding more physical constraints systematically reduces hypothesis complexity and improves generalization behavior.

% 異なる仮定セットに基づくRadmacher Complexityを通した議論は
% Furthermore, while the dependence on intrinsic dimension $d_{\mathcal{V}}$ may evoke comparisons to rank-based arguments, it is essential to recognize a key distinction. In our setting, the design matrix $\mathbf{\Phi}$ is not inherently assumed to be low-rank. Instead, the effective low-rank behavior is an emergent property, arising from the structural constraints imposed by the underlying differential equations and captured by the affine variety $\mathcal{V}$. A significant contribution of our work lies precisely in demonstrating how the behavior of physics-informed models can be understood through the lens of dimensionality reduction concepts analogous to rank, thus providing a novel perspective on the benefits of integrating physical knowledge. Moreover, unlike conventional rank-based analyses, which are often restricted to linear differential equations, our affine variety framework naturally extends to accommodate nonlinear cases, offering a more versatile analytical tool for machine learning and physics-informed modeling. The methodology for computing the intrinsic dimension $d_{\mathcal{V}}$ of the affine variety is detailed in Section~\ref{sec:dim-variety}.

% To qualitatively estimate the impact of the second term, we discuss the case where the generalization capacity is determined by the local size of the hypothesis space induced by  the learning algorithm, such as gradient descent. When using gradient descent for the optimization of~\cref{eq:lin-problem-ball}, the weights \(\bm{w}\) are likely to be trapped in the path-connected component near the initial point because of the local convexity of the loss landscape, as illustrated in~\cref{fig:cover-ofaffineV}. In this situation, the intrinsic size of the hypothesis space can be estimated as follows. The first term of the covering number shown in~\cref{eq:lem:cov} (and consequently in~\cref{eq:thm:main}) remains unchanged because it depends on the dimension. On the other hand, the second term becomes smaller because it focuses on fewer path-connected components. Therefore, we infer that the dimension of the affine variety \(d_{\mathcal{V}}\) in the first term primarily contributes to the generalization capacity, especially when it is determined by the local size of the hypothesis space.

\subsection{Analysis on Linear Operator}
\label{sec:linop}
We discuss the special case where \( \mathscr{D} \) is a linear operator. The second term in \cref{eq:thm:main} vanishes because the Petrovskii-Oleinik-Milnor inequality indicates \( \beta = 1 \). Thus, the minimax risk is \( \mathcal{O}\left(\sqrt{d_{\mathcal{V}} \log (d_{\mathcal{V}}d) / n}\right) \). Furthermore, the affine variety \( \mathcal{V} \) is the solution set of a homogeneous system of linear equations. That is, the affine variety can be written as \( \mathcal{V}(\mathscr{D}, \mathcal{B}, \mathcal{T}) = \{ \bm{w} : \mD \bm{w} = \bm{0} \} \) using the matrix \( \mD \in \mathbb{R}^{K \times d} \) defined by \(\emD_{k, j} \coloneqq \langle \mathscr{D}[\phi_j], \psi_{k} \rangle_{\mu_k}\). The affine variety is a linear subspace of dimension \( d_{\mathcal{V}} = \dim \ker \mD \). From the rank–nullity theorem, \( d_{\mathcal{V}} = d - \operatorname{rank} \mD \), indicating that the higher the rank of the matrix \( \mD \), the better the minimax risk of regression. 

We show that our theory is consistent with existing theories. The effect of incorporating physical structure, represented by linear differential equations, on generalization has been analyzed within the framework of kernel methods by \citet{doumeche2024physics,doumeche2024physics2}. They argued that the physical structure smooths the kernel and reduces the effective dimension, leading to an improvement in the \(\ell_2\) predictive error. We first present the definition of the physics-informed (PI) kernel. 
\begin{definition}[PI kernel~\citep{doumeche2024physics,doumeche2024physics2}]
\label{def:pi-kernel}
Given a basis \( \mathcal{B} = \{ \phi_j \}_{j=1}^d \), trial functions (with a single measure) \( \mathcal{T} = \{ (\psi_k, \mu) \}_{k=1}^K \), and a linear operator \(\mathscr{D}\), the \emph{PI kernel} associated with the affine variety $\mathcal{V}(\mathscr{D}, \mathcal{B}, \mathcal{T}) = \{\bm{w}\in\mathbb{R}^d: \mD\bm{w} = \bm{0}\}$ is defined as:
\begin{align}
\label{eq:def-pi-kernel}
\kappa_{\mM}(x,y) = \left\langle \mM^{-1/2}\bm{\phi}(x),\,\mM^{-1/2}\bm{\phi}(y)\right\rangle_{2},
\end{align}
with
\begin{align*}
\mM(\xi, \nu) \coloneqq \xi \mI + \nu \mD^\top \mT \mD,\quad \mT_{k,k'} = \langle \psi_k,\psi_{k'} \rangle_{\mu}, \quad
\mD_{k,j} = \langle \mathscr{D}[\phi_j],\psi_k\rangle_{\mu}.
\end{align*}
Here, \(\mI\) is the identity matrix, and the matrix $\mT$ is positive semi-definite. The parameters \(\xi,\nu \geq 0\) control the balance between the \(L^2\)-regularization and the constraints derived from the operator \(\mathscr{D}\). 
\end{definition}
% Note that \cref{def:pi-kernel} extends the original definition to more general trial functions. The original PI kernel uses basis functions as trial functions, \ie, \( \mathcal{T} = \mathcal{B} \times \{ \mu \} \). 

\citet{doumeche2024physics2} showed the effective dimension \( d_{\mathrm{eff}}(\xi, \nu) \) of the PI kernel is evaluated above by a computable quantity as follows:
\begin{gather}
    \label{eq:effective-dim-ub}
    d_{\mathrm{eff}}(\xi, \nu) \lesssim \sum_{\alpha \in \sigma\left(\mB\mM^{-1}\mB\right)} \frac{1}{1 + \alpha^{-1}},
\end{gather}
where \( \sigma(\cdot) \) denotes the spectrum (set of eigenvalues) of the given matrix, \(\mB \in \mathbb{R}^{d \times d}\) is the Gram matrix of the basis function, \ie, \(\emB_{j, j'} = \langle \phi_j, \phi_{j'} \rangle_{\mu} \).
Next, we provide an explicit upper bound on the effective dimension of the PI kernel defined using the affine variety:
\begin{restatable}{proposition}{eff}\label{thm:pi-kernel}
    The effective dimension of the PI kernel associated with the affine variety \( \mathcal{V}(\mathscr{D}, \mathcal{B}, \mathcal{T}) = \{ \bm{w} : \mD\bm{w} = \bm{0} \} \) with dimension \( d_{\mathcal{V}} \) is upper bounded by
    \begin{align*}
       d_{\mathrm{eff}}(\xi, \nu) \lesssim  \sum_{j=1}^{d_{\mathcal{V}}} \frac{1}{1 + \xi} + \sum_{j=d_{\mathcal{V}}}^{d} \frac{1}{1 + \xi + \nu \alpha_j} \leq \frac{d}{1 + \xi}.
    \end{align*}
    where \( \{\alpha_j\}_{j=d_{\mathcal{V}}}^d \) denote the positive eigenvalues of the matrix \(\mD^{\top}\mT\mD\).
\end{restatable}

\Cref{thm:pi-kernel} indicates that as the dimension of the affine variety \( d_{\mathcal{V}} = d - \operatorname{rank} \mD \) decreases, the upper bound on the effective dimension of the PI kernel decreases accordingly.
Since the matrix $\mD^{\top}\mT\mD$ is positive semi-definite, all eigenvalues satisfy $\alpha_j \geq 0$. The intrinsic dimension $d_{\mathcal{V}}$ corresponds precisely to the number of zero eigenvalues ($\alpha_j = 0$). The terms in the second sum ($j > d_{\mathcal{V}}$) involve strictly positive eigenvalues $\alpha_j > 0$. Given that $\nu > 0$, we have $\frac{1}{1 + \xi + \nu\alpha_j} < \frac{1}{1 + \xi}$. Thus, when the intrinsic dimension $d_{\mathcal{V}}$ decreases, the number of terms in the first sum (with the larger value $1/(1 + \xi)$) decreases, while the number of terms in the second sum (with smaller values $1/(1 + \xi + \nu\alpha_j)$) increases. This shift towards smaller-valued terms leads to an overall reduction in the complexity bound.

Consequently, our theoretical results align with the existing PI kernel theory~\citep{doumeche2024physics,doumeche2024physics2}. The PI kernel framework from the previous literature quantifies the complexity of the hypothesis space through the entire spectrum of the matrix $\mD$ combined with the base kernel $\langle \bm{\phi}(x), \bm{\phi}(y) \rangle_{2}$, restricting the analysis primarily to linear target operators $\mathscr{D}$. In contrast, our approach allows for analysis of linear and non-linear operators by focusing solely on the intrinsic dimension $d_{\mathcal{V}}$ (the count of zero eigenvalues), rather than analyzing the entire eigenvalue spectrum.


\section{On the Dimension of an Affine Variety}
\label{sec:dim-variety}
In general, the dimension of the affine variety \( V = \{ \bm{w} \in \mathbb{R}^d: p_k(\bm{w}) = 0, \forall k = 1,\ \ldots,\ K \} \) defined by polynomials $\{ p_k \}_{k=1}^K$ has many equivalent definitions. In particular, the following statements are all equivalent. 
\begin{restatable}{definition}{maxChain}\label{def:max-chain}
    The maximal length of the chains \( V_0 \subset V_1 \subset \ldots \subset V_{d_V} \) of non-empty subvarieties of \( V \).
\end{restatable}
\begin{restatable}{definition}{hilbertSeries}\label{def:hilbert-series} The degree of the denominator of the Hilbert series of the affine variety \(V\).
\end{restatable}
\begin{restatable}{definition}{tangentSpace}\label{def:tangent-space} The maximal dimension of the tangent vector spaces at the non-singular points \( U \subseteq V \subset \mathbb{R}^d \) of the variety, \ie, $d_V = \max_{\bm{w} \in U} d - \operatorname{rank}
    \begin{bmatrix} 
        \nabla p_1(\bm{w}) &
        \cdots &
        \nabla p_K(\bm{w}) 
    \end{bmatrix}^{\top}$.
\end{restatable}
%\end{description}

Although \cref{def:max-chain} clearly indicates that the dimension represents the complexity of the set \( V \), it is difficult to calculate the dimension according to this definition. \Cref{def:hilbert-series} shows that the dimension represents the algebraic complexity of the polynomial ring. \Cref{def:tangent-space} characterizes the dimension based on the local structure of the affine variety, making it suitable for numerical calculation as discussed in~\cref{sec:dv-cal}. It generalizes the rank-nullity theorem \( d_{V} = d - \operatorname{rank} \mD \) in the linear case, as mentioned in~\cref{sec:linop}. The details of the concepts associated with these definitions are given in~\cref{app:affine-variety}. 

\subsection{Lower Bound}
We demonstrate that the dimension \( d_{\mathcal{V}} \) of the affine variety can be characterized by the linear part of the operator \( \mathscr{D} \). 
\begin{restatable}{proposition}{lowerbound}
\label{thm:lb-dim}
Suppose the operator \( \mathscr{D} \) can be decomposed as \( \mathscr{D} = \mathscr{L} + \mathscr{F} \), where \( \mathscr{L} \) is a nonzero linear differential operator and \( \mathscr{F} \) is a nonlinear operator. Then, we have $d_{\mathcal{V}(\mathscr{L})} \leq d_{\mathcal{V}(\mathscr{D})}$.
\end{restatable}

Combining the result of~\cref{thm:lb-dim} with~\cref{thm:main} suggests that the nonlinear part \( \mathscr{F} \) of the operator increases the affine variety dimension, which has a negative effect on generalization. Furthermore, the dimension of the affine variety associated with the linear part \( \mathscr{L} \) can be easily calculated by the rank of the matrix. Therefore, the lower bound of the dimension of the affine variety associated with the nonlinear operator \( \mathscr{D} \) can be easily determined, allowing us to estimate the minimum required amount of data \( n \).

\subsection{Numerical Calculation Method}
\label{sec:dv-cal}
According to \Cref{def:hilbert-series}, the dimension of an affine variety is typically obtained by calculating the degree of the denominator of the Hilbert series, by using Gröbner bases. 
However, the worst-case time complexity of Buchberger's algorithm~\citep{buchberger1976theoretical}, the standard method for computing Gröbner bases, is double exponential in the number of variables \(d\). Therefore, on the basis of~\ref{def:tangent-space}, we approximate \(d_{\mathcal{V}}\) by sampling \(\bm{w}^*_{1},\ \ldots,\ \bm{w}^*_{N}\) from the affine variety \(\mathcal{V}\) with a suitable distribution and then computing \(\max_{\bm{w}^* \in \{\bm{w}^*_1,\ldots,\bm{w}^*_N\}} d-\mathrm{rank}\left(\nabla^{\top} [p_1(\bm{w}^*),\ \ldots,\ p_{K}(\bm{w}^*)]^{\top}\right)\). 
When the operator \(\mathscr{D}\) is nonlinear, we perform simulations with various boundary conditions and project the obtained solutions onto the basis \(\mathcal{B}\) to sample \(\bm{w}^* \in \mathcal{V}\). For linear operators, the dimension does not depend on the particular weight \(\bm{w}\), and the rank of the matrix \(\mD\) discussed in~\cref{sec:linop} precisely determines \(d_{\mathcal{V}}\). 
Assuming the use of standard rank computation algorithms, the computational complexity of this numerical approach is \(\mathcal{O}(N \cdot \min(K, d) K d)\) for the nonlinear case, and \(\mathcal{O}(\min(K, d) K d)\) for the linear case. This complexity is practical and feasible for most scenarios considered in our setting.

\section{Experiments}
\label{sec:exp}
To evaluate the generalization performance of physics-informed linear regression (PILR) compared to ridge regression (RR) using basis functions \(\mathcal{B}\), we conducted experiments on representative differential equations. We varied the data size \(n\) and parameter count \(d\), and report test MSE (mean ± standard deviation) across 10 random initial or boundary conditions. Experimental details are provided in~\cref{app:exp_detail}.

When the operator \(\mathscr{D}\) is linear, PILR approximates the solution to~\cref{eq:lin-problem-ball} as~\citep{doumeche2024physics2}:
\begin{align*}
\widehat{\bm{w}} = (\bm{\Phi}^{\top}\bm{\Phi} + n\mM)^{-1} \bm{\Phi}^{\top} \bm{y},
\end{align*}
where \(\mM\) depends on hyperparameters \(\xi\) and \(\nu\) (see~\cref{eq:def-pi-kernel}); setting \(\nu = 0\) yields RR.

For nonlinear equations, we train models by minimizing a soft-constrained loss using the Adam optimizer. Hyperparameters \(\xi\) and \(\nu\) are tuned via validation MSE. 

\subsection{Learning Strong Solutions}
\label{sec:strong-sol}
\begin{figure*}[tb]
    \captionsetup[subfigure]{justification=centering}
    \centering
    \subcaptionbox{Harmonic Oscillator\label{fig:ho-perf}}[0.35\textwidth]{
    \includegraphics[width=0.35\textwidth]{fig/results_osc.png}
    }
    \centering
    \subcaptionbox{Diffusion Eq.\label{fig:diff-perf}}[0.35\textwidth]{
    \includegraphics[width=0.35\textwidth]{fig/results_diff.png}
    }%
    \centering
    \subcaptionbox{Predictions of Harmonic Oscillator \label{fig:ho-vis}}[0.29\textwidth]{
    \includegraphics[width=0.29\textwidth]{fig/Oscillator_vis.png}
    }
    \caption{Experimental results for the strong solutions. (a, b) Test MSE (log scale) vs. number of parameters for the harmonic oscillator (a) and diffusion equation (b). The plots compare RR and PILR for three different data sizes \(n\), showing the mean and standard deviation across 10 initializations. (c) Predictions of harmonic oscillator using a 33-parameter model trained on 20 samples: RR and PILR, with training data points indicated.}
    \label{fig:strong-sol}
\vskip -0.2in
\end{figure*}
In this section, we investigate the strong solutions of the classical harmonic oscillator and the diffusion equation with periodic boundary conditions, by employing the Dirac measure, which corresponds to the collocation method used in PINNs. The solutions to these equations can be obtained analytically. Through these straightforward examples, we demonstrate both analytically and numerically that the generalization performance is determined by the dimension of the affine variety. 

\paragraph{Harmonic Oscillator} The initial value problem of a harmonic oscillator \(\mathscr{D}[y]=0\) with a spring constant \( k_s \) and mass \( m_s \) in the domain \( \Omega = [0, T] \) is given by:
\begin{align}
    \label{eq:osc-eq}
    \mathscr{D}[y] = \frac{\mathrm{d}^2}{\mathrm{d}t^2}y + \frac{k_s}{m_s}y,\quad y(0)=y_0,\quad \frac{\mathrm{d}}{\mathrm{d}t}y(0)=v_0,
\end{align}
where \(y_0\) and \(v_0\) are the initial position and velocity, respectively. The solution to the initial value problem is analytically given by \( y(t) = y_0 \cos(\omega t) + \frac{v_0}{\omega} \sin(\omega t), \ \omega = \sqrt{k_s/m_s} \). 
The settings for the basis and the trial functions with the measure \(\phi_j \in \mathcal{B},\ (\psi_k, \mu_k) \in \mathcal{T}\) of indices \(1 \leq j \leq d_t\) and \(1 \leq k \leq K_t \) are as follows:
\begin{align}
    \label{eq:basis_test_osc}
    \phi_1(x) = 1,\ \phi_{2j}(x) = \cos\left(\omega_j x\right),\ 
    \phi_{2j+1}(x) = \sin\left(\omega_j x\right),\
    \psi_k(x) = 1,\ \mu_k = \delta_{x_k},
\end{align}
where \(\omega_j \coloneqq \frac{j\pi}{T}\) is the \(j\)-th frequency and \(\delta_{x_k}\) is the Dirac measure centered at the point \(x_k \in \Omega\), which is uniformly sampled from data. 

Then, the dimension of the affine variety is \( d_{\mathcal{V}} = 2\), representing the essential degrees of freedom of the solution. \Cref{fig:ho-perf} supports our theory experimentally. For RR, the generalization performance degrades as the number of parameters \( d = 2d_t + 1 \) increases due to overfitting, as shown in~\cref{fig:ho-vis}. In contrast, for PILR, the performance remains stable regardless of the number of parameters \( d \) owing to the lower dimension of the affine variety \( d_{\mathcal{V}} = 2 \).

\paragraph{Diffusion Equation} The initial value problem for the one-dimensional diffusion equation \(\mathscr{D}[u]=0\) with diffusion coefficient \( c \) and periodic boundary conditions is given by:
\begin{align}
\label{eq:diffusion-eq}
\begin{cases}
 \frac{\partial u}{\partial t} -c \frac{\partial^2 u}{\partial x^2}= 0,
 & (x,t)\in[-\Xi,\Xi]\times[0,T] \\
 u(x,0)=u_0(x),
 & x\in[-\Xi,\Xi] \\
 u(-\Xi,t) = u(\Xi,t),\;
 \frac{\partial u}{\partial x}(-\Xi,t) = \frac{\partial u}{\partial x}(\Xi,t),
 & t\in[0,T]
\end{cases}
\end{align}

We define the basis functions $\phi \in \mathcal{B}$ and the test functions with measures $(\psi, \mu) \in \mathcal{T}$ as follows:
\begin{align}
\label{eq:basis_test_diff}
\phi_{2j, j'} = \cos(\omega_j x) e^{-c\omega_{j'}^2 t},\ \phi_{2j+1, j'} = \sin(\omega_j x) e^{-c \omega_{j'}^2t},\quad
\psi_{k, k'} = 1,\ \ \mu_{k, k'}= \delta_{(t_k, x_{k'})},
\end{align}
where the frequency is $\omega_j = j\pi/\Xi$. The indices are in the ranges $0 \leq j \leq d_x$, $0 \leq j' \leq d_t$, $1 \leq k \leq K_t$, and $1 \leq k' \leq K_x$.

The analytical solution is expressed as a linear combination of the above basis functions. The number of bases is \(d = 2d_xd_t + 1 \), while the dimension of an affine variety is given by \( d_{\mathcal{V}} = 2 \min(d_x, d_t) + 1\). \Cref{fig:diff-perf} shows the results when we set \(\alpha = 1.0\), \(j_{\mathrm{max}} = 1\), \(d_t = 2\), and vary \(d_x\). The results indicate that the generalization performance of PILR does not deteriorate as \(d_x\) increases, in contrast to RR. 

\subsection{Learning Weak Solutions}
\label{sec:weak-sol}
\begin{figure*}[tb]
    \captionsetup[subfigure]{justification=centering}
    \centering
    \subcaptionbox{Harmonic Oscillator\label{fig:ho-perf-weak}}[0.35\textwidth]{
    \includegraphics[width=0.35\textwidth]{fig/results_osc_weak.png}
    }
    \centering
    \subcaptionbox{Diffusion Eq.\label{fig:diff-perf-weak}}[0.35\textwidth]{
    \includegraphics[width=0.35\textwidth]{fig/results_diff_weak.png}
    }%
    \centering
    \caption{Experimental results for the weak solutions. (a, b) Test MSE (log scale) vs. number of parameters for the harmonic oscillator (a) and diffusion equation (b). The plots compare RR and PILR for three different data sizes \(n\), showing the mean and standard deviation across 10 initializations.}
    \label{fig:weak-sol}
\vskip -0.2in
\end{figure*}
In this section, we investigate weak solutions for the harmonic oscillator and the diffusion equation, employing a variational framework with Borel measures. The governing equations and basis functions are identical to those in \cref{sec:strong-sol}.

\paragraph{Harmonic Oscillator}
We define the trial functions $\psi_k$ for $1 \leq k \leq K_t$ as:
\begin{equation}
    \psi_1(x) = 1, \quad \psi_{2k-1}(x) = \cos(\omega_k x), \quad \psi_{2k}(x) = \sin(\omega_k x),
\end{equation}
where $\omega_k \coloneqq \frac{k\pi}{T}$ is the frequency. The associated measure $\mu_k$ is the Lebesgue measure on $\Omega = [0, T]$.

The dimension of the affine variety remains $d_{\mathcal{V}} = 2$, consistent with the strong solutions. Experimental results in \cref{fig:ho-perf-weak} confirm this. The performance of PILR is stable and independent of the number of basis functions, unlike RR, which shows performance degradation as model complexity increases.

\paragraph{Diffusion Equation}
The trial functions $\psi_{k, k'}$ combine a piecewise constant basis in time and a Fourier basis in space. For indices $1 \leq k \leq K_t$ and $1 \leq k' \leq K_x$, they are:
\begin{equation}
    \psi_{k, 2k'-1}(x,t) = 1_{[t_k, t_{k+1}]}(t) \cos(\omega_{k'} x), \quad \psi_{k, 2k'}(x,t) = 1_{[t_k, t_{k+1}]}(t) \sin(\omega_{k'} x),
\end{equation}
where $1_{[t_k, t_{k+1}]}(t)$ is the indicator function for the time interval, defined as
\begin{equation}
    1_{[t_k, t_{k+1}]}(t) = 
    \begin{cases} 
        1 & \text{if } t \in [t_k, t_{k+1}) \\ 
        0 & \text{otherwise} 
    \end{cases},
\end{equation}
and $\omega_{k'} = \frac{k'\pi}{\Xi}$. The associated measure is the Lebesgue measure on $[-\Xi, \Xi] \times [0, T]$.

The dimension of the affine variety, $d_{\mathcal{V}} = 2 \min(d_x, d_t) + 1$, is identical to the strong solution case. The results in \cref{fig:diff-perf-weak} show that PILR's generalization performance remains robust as the number of spatial basis functions $d_x$ increases, demonstrating its advantage over RR.

\subsection{Learning Numerical Solutions}
\label{sec:numel-sol}
In this section, we learn approximate solutions using numerical methods that use finite difference for four equations. In this setting, we consider the affine variety of the difference equation \(\mathscr{D}_{\bm{h}}\) and the base functions \(\mathcal{B}_{\bm{h}}\) and the trial functions with the measure \(\mathcal{T}_{\bm{h}}\) corresponding to the numerical method with step size \(\bm{h}\). We first validate our theory using linear and nonlinear Bernoulli equations discretized by the explicit Euler method. 
\paragraph{Discrete Bernoulli Equation}
We discretize the Bernoulli equation on the interval $\Omega=[0,T]$ with uniform step size~$h$:
\begin{align*}
  \mathscr{D}_h[y]
  \;=\;\frac{y_{\tau+1}-y_{\tau}}{h} \;+\; P\,y_{\tau}\;-\;Q\,y_{\tau}^{\rho}
  \;=\;0, 
  \quad \tau=0,\dots,n_t-1,\quad n_t=\frac{T}{h},
\end{align*}
where \( y_{\tau} = y(t_{\tau}) \). We consider two parameter regimes \( (P, Q, \rho) \) set to \((1.0, 0.0, 0.0)\) for the linear case and to \((1.0, 0.5, 2)\) for the non-linear case.
The initial value~$y_0$ is sampled from $\mathcal{N}(0,1)$, and the reference solution is calculated explicitly by Euler.  
Further details on the choice of $n_t$, basis/trial functions, measure $(\psi_\tau,\mu_\tau)$, and implementation are given in~\cref{app:exp-numel-sol}. 

\paragraph{Discrete Diffusion Equation}
We discretize the one‐dimensional diffusion equation over 
\(\Omega=[-\Xi,\Xi]\times[0,T]\) with step sizes \(\bm h=(h_x,h_t)\) and diffusion coefficient \(c(u)\):
\begin{align*}
  \mathscr{D}_{\bm h}[u]
  &= \frac{u_j^{\tau+1}-u_j^{\tau}}{h_t}
     -c(u_j^{\tau})\,\frac{u_{j+1}^{\tau}-2\,u_{j}^{\tau}+u_{j-1}^{\tau}}{h_x}
  \;=\;0,
  \quad j = 1, \dots, n_x,\;\tau = 1, \dots, n_t,
\end{align*}
where \(u_j^{\tau}=u(x_j,t_{\tau})\).  
We consider two cases:  \( c(u) = 1.0 \) for the linear case and \( c(u) = 0.1/(1 + u^2) \) for the nonlinear case.
Periodic boundary conditions are imposed in \(x\).  
More details on the grid, the basis / trial functions, and the numerical setup are given in~\cref{app:exp-numel-sol}. 

\Cref{tab:euler-bernoulli,tab:fdm-diffusion} show that PILR achieves a higher performance than RR for large values of \(d\). While the dimension \(d_{\mathcal{V}}\) is independent of the time discretization step size in the Euler method, it depends on the spatial discretization step size in the FDM. We include supplementary experiments in~\cref{app:additional_exp}, where we fix the ambient dimension \(d\) and vary the size of the trial‐function set \(\mathcal{T}\). 

\begin{table*}[tb]
\caption{Experimental results for the discrete linear and nonlinear Bernoulli equations approximated by the explicit Euler method. The settings include various step sizes \(h\). The number of parameters (basis) \(d\), and the calculated dimension of the affine variety \(d_{\mathcal{V}}\).}
\label{tab:euler-bernoulli}
\begin{center}
\begin{tabular}{cc|cc|cc}
\toprule
\multirow{2}{*}{Settings} & \(\mathscr{D}_h\) & \multicolumn{2}{c|}{Linear Bernoulli eq.} & \multicolumn{2}{c}{Nonlinear Bernoulli eq.} \\ 
& \( h \) & \( 1/100 \) & \( 1/200 \) & \( 1/100 \) & \( 1/200 \)  \\
\multirow{2}{*}{Dimensions} & \( d \) & 100 & 200 & 100 & 200 \\
& \(d_{\mathcal{V}}\) & 1 & 1 & 1 & 1  \\ \midrule 
\multirow{2}{*}{Test MSE} & RR & \(0.48 \pm 0.32\) & \(0.63 \pm 0.43 \) & \(0.60 \pm 0.41\) & \(0.72 \pm 0.49\) \\
& PILR & \(0.012 \pm 0.0025\) & \(0.011 \pm 0.0013\)  & \(0.013 \pm 0.0024\) & \(0.013 \pm 0.0018\)  \\ \bottomrule
\end{tabular}
\end{center}
\vskip -0.1in
\end{table*}

\begin{table*}[tb]
\caption{Experimental results for the discrete linear and nonlinear diffusion equations approximated by the FDM. The settings include various step sizes \(\bm{h} = (h_t, h_x)\). The number of parameters (basis) \(d\), and the calculated dimension of the affine variety \(d_{\mathcal{V}}\).}
\label{tab:fdm-diffusion}
\begin{center}
\begin{tabular}{cc|cc|cc}
\toprule
\multirow{2}{*}{Settings} & \(\mathscr{D}_h\) 
  & \multicolumn{2}{c|}{Linear diffusion eq.} 
  & \multicolumn{2}{c}{Nonlinear diffusion eq.} \\ 
 & \((h_t, h_x)\) 
  & \(\left(1/400, 2/10\right)\) 
  & \(\left(1/400, 2/20\right)\) 
  & \(\left(1/200, 2/10\right)\) 
  & \(\left(1/200, 2/20\right)\) \\ 
\multirow{2}{*}{Dimensions} 
 & \(d\) 
  & 4010 & 8020 
  & 2010 & 4020 \\ 
 & \(d_{\mathcal{V}}\) 
  & 10   & 20   
  & 10   & 20   \\ 
\midrule 
\multirow{2}{*}{Test MSE} 
 & RR  
  & \(2.21 \pm 0.56\) & \(2.14 \pm 0.57\) 
  & \(1.12 \pm 0.40\) & \(1.11 \pm 0.40\) \\ 
 & PILR
  & \(1.13 \pm 0.30\) & \(0.79 \pm 0.16\) 
  & \(0.26 \pm 0.11\) & \(0.22 \pm 0.10\) \\ 
\bottomrule
\end{tabular}
\end{center}
\vskip -0.1in
\end{table*}

\subsection{Impact of Basis Misspecification on Generalization}

This section considers a practical scenario where the basis functions are misspecified, a situation that can occur during manual design or through random selection, as in an Extreme Learning Machine (ELM)~\citep{huang2006extreme}. Any such misspecification can degrade performance by increasing the \textbf{approximation error}. As detailed in~\cref{app:analysis setting}, the total error is composed of this approximation error and an estimation error. While our theory demonstrates that physical constraints can reduce the estimation error, the overall model performance is limited by the magnitude of the approximation error. 

To demonstrate this effect, we conducted an experiment on the Harmonic Oscillator, intentionally omitting the known analytical frequency from the basis functions. Other experimental settings were identical to those in Section 5.1. With 10 data points, the performance was exceptionally poor. For a basis size of $d=17$, the test MSE was approximately $1.435 \pm 0.646$, of which the approximation error constituted nearly the entire amount at $1.430$. Increasing the basis size to $d=33$ had a negligible effect; the test MSE remained high at $1.434$ as the approximation error was unchanged.

This result clearly shows the total error being dominated by the approximation error. It underscores a prerequisite for our theory: the improvement in generalization from physics-informed constraints is achieved only when the model possesses sufficient expressive capacity to represent the true solution. 
% The systematic selection of an appropriate basis to ensure this capacity is a non-trivial model selection problem, which remains an important direction for future research.

\section{Conclusion}
This study introduces a framework for analyzing physics-informed models through the lens of affine varieties induced by the governing differential equations. We establish that generalization performance is governed by the dimension of this variety, rather than the number of model parameters, a finding that unifies existing theories for linear equations. We further provide a method for calculating this dimension and present experimental validation confirming that this intrinsic dimension effectively mitigates overfitting in highly parameterized settings. Although our analysis centers on linear regression models, the proposed geometric framework is broadly applicable to both linear and nonlinear differential equations, as our experiments demonstrate. This work offers a foundational, geometric interpretation of generalization that establishes a promising, though challenging, direction for future theory-guided model selection, such as the optimal choice of basis and trial functions. Future work includes the extension and validation of our framework for other architectures, such as NN and ELM. The framework can also be extended to differential equations with unknown parameters by analyzing an augmented parameter space.

% Adapting our analysis to conservation laws \citep{jagtap2020conservative, hu2022extended} or Lie symmetries \citep{akhound2024lie, dalton2024physics} is a promising direction. 

% In the unusual situation where you want a paper to appear in the
% references without citing it in the main text, use \nocite
\bibliography{example_paper}
\bibliographystyle{plainnat}

\newpage
\appendix
\section{Notation}
\label{app:notation}
\begin{table}[!htbp]
\centering
\begin{tabular}{ll}
\textbf{Symbol} & \textbf{Description} \\
\midrule
\multicolumn{2}{l}{\vspace{0.02in} \textbf{Data}} \\
$f^*$ & True function to be learned \\
$\Omega \subseteq \mathbb{R}^m$ & Input domain \\
$m$ & Input dimension \\
$n$ & Number of observations \\
$(x_i, y_i)$ &  $i$-th observation ($x_i$: input, $y_i$: output) \\
$\boldsymbol{y}$ & Target vector $[y_1, \dots, y_n]^\top$ \\
$\sigma^2$ & Noise variance \\
$\varepsilon_i$ & Normally distributed noise following $\mathcal{N}(0, \sigma^2)$ \\
\midrule
\multicolumn{2}{l}{\vspace{0.02in} \textbf{Affine Variety and Variables}} \\
$\mathscr{D}$ & Differential operator \\
$\mathcal{B} = \{\phi_j\}_{j=1}^d$ & Basis functions \\
$\phi_j$ & j-th basis function from $\mathcal{B}$ \\
$\bm{\phi}(x) \in \mathbb{R}^d$ & Basis vector at $x$ \\
$\mathbf{\Phi} \in \mathbb{R}^{n \times d}$ & Design matrix (i-th row is $\boldsymbol{\phi}(x_i)^\top$) \\
$\mathcal{T} = \{(\psi_k, \mu_k)\}_{k=1}^K$ & Finite collection of trial function and measure pairs. \\
$d$ & Number of basis functions $|\mathcal{B}|$ (ambient dimension) \\
$K$ & Number of trial functions $|\mathcal{T}|$ \\
$\mathcal{V}(\mathscr{D}, \mathcal{B}, \mathcal{T})$ & Affine variety defined by $\mathscr{D}, \mathcal{B}, \mathcal{T}$ (set of weight vectors) \\
$d_{\mathcal{V}}$ & Dimension of the affine variety $\mathcal{V}$ \\
$\bm{w}$, $\hat{\bm{w}}$, $\bm{w}^*$ & Weight vectors (learnable, estimated, optimal) \\
$\mathbb{B}_2(R)$ & $\ell_2$ ball of radius $R$ \\
$\mathcal{V}_R$ & Affine variety constrained by the $\ell_2$-ball ($\mathcal{V} \cap \mathbb{B}_2(R)$) \\
$\lambda_n$ & $L^2$ regularization parameter \\
$\|\cdot\|_2$, $\|\cdot\|$ & Vector $\ell_2$ norm, function $L^2$ norm w.r.t. Borel measure \\
$\langle \cdot, \cdot \rangle_{2},\ \langle \cdot, \cdot \rangle_{\mu}$ & Vector Euclidean inner product, function inner product w.r.t. measure $\mu$ \\
\midrule
\multicolumn{2}{l}{\textbf{\vspace{0.02in} Geometric / Complexity Measures}} \\
$V \subseteq \mathbb{R}^d$ & General affine variety \\
$(\beta, d_V)$-regular set & Regularity condition \\
$d_V$ & Dimension of $V$ \\
$\operatorname{codim}(L)$ & Codimension $d - d_L$ \\
$\beta$ & Upper bound on connected components of intersections \\
$\mathcal{N}(V, \varepsilon, \|\cdot\|_2)$ & $\varepsilon$-covering number w.r.t. $\ell_2$ norm $\|\cdot\|_2$.\\
\midrule
\multicolumn{2}{l}{\vspace{0.02in}\textbf{Analysis Constants}} \\
$M$ & Upper bound constant for the basis functions, such that $\|\boldsymbol{\phi}(x)\|_2 \leq M$ \\
$\eta$ & Upper bound constant for the lower eigenvalue of design matrix $\mathbf{\Phi}$.  \\
$\Gamma$ & Stability constant of the estimator \\
$\delta$ & Probability parameter \\
\midrule
\multicolumn{2}{l}{\vspace{0.02in}\textbf{Linear Operators / PI Kernel}} \\
$\mD, \emD_{k,j}$ & Constraint matrix when the operator $\mathscr{D}$ is linear; $\emD_{k,j}$ is its entry \\
$\mathscr{L}, \mathscr{F}$ & Linear, nonlinear parts of $\mathscr{D}$ \\
$d_{\mathcal{V}(\mathscr{L})}$, $d_{\mathcal{V}(\mathscr{D})}$ & Dimensions under $\mathscr{L}$, $\mathscr{D}$ \\
$\kappa_{\mM}(x,y)$ & PI kernel defined with regularization matrix $\mM$ \\
$\xi, \nu$ & Hyperparameters of the PI kernel (controlling the balance)  \\
$\mB$, $\emB_{j,j'}$ & Basis Gram matrix, its entry \\
$\mT$, $\emT_{k,k'}$ & Trial Gram matrix, its entry \\
$d_{\mathrm{eff}}(\xi, \nu)$ & Effective dimension of the PI kernel \\
$\sigma(\cdot),\ \alpha, \alpha_j$ & Spectrum of a matrix and its entry (eigenvalues)  \\
\midrule
\multicolumn{2}{l}{\vspace{0.02in}\textbf{Dimension Calculation}} \\
$p_k$ & Defining polynomial \\
$N$ & Number of samples $(\bm{w}^*_1,\ \ldots,\ \bm{w}^*_N)$ \\
\bottomrule
\end{tabular}
\end{table}

\newpage
\section{Mathematical Background on Affine Varieties}
\label{app:affine-variety}
\begin{figure}[tb]
    \centering
    \includegraphics[width=0.4\linewidth]{fig/V2.png}
    \caption{Illustration of the construction of the \(\varepsilon\)-covering of the affine variety \(\mathcal{V} \subseteq \mathbb{R}^2\) and the associated loss landscape. The black curve represents a \((\beta, d_{\mathcal{V}})\) regular affine variety with dimension \(d_{\mathcal{V}} = 1\). The color gradients depict the loss landscape \(\mathcal{L}(\bm{w}) \coloneqq \sum_{k=1}^K \|p_k(\bm{w})\|_2^2\) of the equations defining \(\mathcal{V} = \{ \bm{w} : p_k(\bm{w}) = 0,\ \forall k =1,\ \ldots,\ K\}\). The blue dotted line represents a \(\ell_2\) ball of radius \(R\). The affine variety constrained with the \(\ell_2\) ball is covered by \(\varepsilon\)-balls centered at the intersections of \(\mathcal{V}\) with four given subspaces \(\{L_s\}_{s=1}^4\), shown as red points. The upper bound on the number of intersections of every subspace with the variety is \(\beta\), while the actual maximum number is \(5\) formed by the subspace \(L_{\star}\) (the yellow dotted line). The loss landscape of the equations is zero on \(\mathcal{V}\) and locally convex around the points in \(\mathcal{V}\).}
    \label{fig:cover-ofaffineV}
\end{figure}
In this section, we provide a formal definition of several concepts related to affine varieties and review the definition of the dimension of an affine variety, as briefly described in~\cref{sec:dim-variety}.

An affine variety is a fundamental concept in algebraic geometry. It is a subset of an affine space, defined as the solution set to a system of polynomial equations. Let \( \mathbb{K}[\bm{w}] \) denote the set of polynomials in the variables \(\bm{w} = (w_1, \ldots, w_d) \in \mathbb{K}^d\) over a field \(\mathbb{K}\) (often \(\mathbb{R}\) or \(\mathbb{C}\)). An affine variety \( V(p_1, \ldots, p_K) \subseteq \mathbb{K}^d\) defined by the polynomials \(p_1, \ldots, p_K \in \mathbb{K}[\bm{w}]\) is given by:
\begin{gather*}
V(p_1, \ldots, p_K) \coloneqq \left\{ \bm{w} \in \mathbb{K}^d : p_k(\bm{w}) = 0,\ \forall k = 1,\ \ldots,\ K \right\}.
\end{gather*}

The geometry of an affine variety is determined by the set of all polynomials that "vanish" on \(V\), i.e., those that become zero for every point in \(V\). This set is called the ideal of the affine variety, denoted \(I(V)\), and is defined as follows:
\begin{gather*}
I(V) := \left\{ p \in \mathbb{K}[\bm{w}] : p(\bm{w}) = 0,\ \forall \bm{w} \in V \right\}.
\end{gather*}
The generating polynomial set \(\{ p_k \}_{k=1}^K\) of the affine variety \(V\) is a subset of the ideal \(I(V)\).

The coordinate ring over \(V\), denoted \(\mathbb{K}[V]\), is introduced to identify polynomials that yield the same values on the variety \(V\). Specifically, \(\mathbb{K}[V]\) is defined as the quotient of the polynomial ring \(\mathbb{K}[\bm{w}]\) by the ideal \(I(V)\), i.e., \(\mathbb{K}[\bm{w}] / I(V)\).
In the coordinate ring \(\mathbb{K}[V] = \mathbb{K}[\bm{w}] / I(V)\), the difference between \(p\) and \(q\) vanishes on \(V\), i.e., \(p(\bm{w}) = q(\bm{w})\) for all \(\bm{w} \in V\), or equivalently \(p - q \in I(V)\). Thus, \(p\) and \(q\) are considered the same element. From another viewpoint, the coordinate ring \(\mathbb{K}[V]\) can be considered as a set of polynomials not included in the ideal \(I(V)\).

Based on the above definitions, we review the definition of the dimension \(d_{V}\) of the affine variety in~\cref{app:dimension_of_affinev} and the regularity in~\cref{app:regularity_of_affinev}.

\subsection{Dimension of Affine Varieties}
\label{app:dimension_of_affinev}
\subsubsection{Geometric View}
Considering the affine variety \(V\) as an affine space, we can naturally define a subvariety as an "subset" of the variety that also satisfies polynomial equations. Let \(q_1, \ldots, q_S\) be polynomials in a ring.
Define $\langle q_1, \ldots, q_S\rangle$ as the smallest ideal generated by $q_1, \ldots, q_S$; that is, $\langle q_1, \ldots, q_S\rangle$ consists of all finite sums of the form $\sum_{i=1}^S r_i q_i$ where each $r_i$ is in the ring: $\langle q_1, \ldots, q_S\rangle=\{\sum_{i=1}^S r_iq_i\}$.
A subvariety $U$ of $V$ is defined as the zero set of a subset ideal $\langle q_1, \ldots, q_S\rangle \subseteq \mathbb{K}[\bm{w}]/I(V)$ given by:
\begin{gather*}
U \coloneqq \left\{ \bm{w} \in \mathbb{K}^d : q_s(\bm{w}) = 0,\ \forall q_s \in \langle q_1, \ldots, q_S\rangle\right\}.
\end{gather*}

By using the concept of subvarieties, the dimension of an affine variety is defined as follows:
\maxChain*
This definition intuitively represents the size of \(V\) by the maximal length of an increasing sequence of subspaces. If the generating polynomials \(\{ p_k \}_{k=1}^K\) are all linear, the dimension of \(V\) is defined as the maximal length of an increasing sequence of linear subspaces within \(V\), which corresponds to the dimension of \(V\) as a linear space.

When we focus on the local structure, the following equivalent definition is obtained:
\tangentSpace*
From this definition, we can see that the dimension \(d_V\) is a global quantity that summarizes the local linearized structure of the affine variety \(V\) at a point.

For example, let $\mathbb{K}=\mathbb{R}$ and $V \subset \mathbb{R}^3$ be the plane: $V=\{(x, y, z): x+y-z=0\}$. A chain of subvarieties within $V$ is $V_0 \subset V_1 \subset V_2$, where $V_0=\{(0,0,0)\}$ (a point, $0$-dimensional), $V_1=$ $\{(t, 0, t): t \in \mathbb{R}\}$ (a line, $1$-dimensional), and $V_2=V$ itself (the plane, $2$-dimensional). 
The maximal length of the nested subvarieties is two, \ie, $d_V=2$, which means that a plane has two degrees of freedom. 

\subsubsection{Algebraic View}
The structure of an affine variety is determined by the ideal \(I(V)\). Intuitively, the larger \(I(V)\) is, the more polynomial constraints there are, which means that \(V\) becomes smaller, and consequently, the coordinate ring \(\mathbb{K}[V]\) also becomes smaller. From this perspective, it is natural to expect a deep connection between the dimension of the coordinate ring \(\mathbb{K}[V]\) (and similarly the ideal \(I(V)\)) and the dimension of the affine variety \(V\). 

To explore this connection, we first discuss the dimension of the coordinate ring \(\mathbb{K}[V]\) using Krull dimension. The ideal \(\mathfrak{p} \subset \mathcal{R}\) in a polynomial ring \(\mathcal{R}\) is prime if \(\forall a,b \in \mathcal{R},\ ab \in \mathfrak{p} \Rightarrow a \in \mathfrak{p} \ \text{or} \ b \in \mathfrak{p} \). The definition of the dimension of the affine variety through the Krull dimension is shown below.
\begin{definition}
    The Krull dimension of the coordinate ring \(\mathbb{K}[V]\): The maximum length \(d\) of the chain of prime ideals \(\mathfrak{p}_0 \subset \mathfrak{p}_1 \subset \dots \subset \mathfrak{p}_d \) in the coordinate ring \(\mathbb{K}[V]\). 
\end{definition}
This definition signifies that the dimension of an affine variety is characterized in the world of polynomial sets by the maximal length of an increasing chain of ``subsets'' within the coordinate ring, corresponding to~\cref{def:max-chain} from a geometric perspective. 

In contrast, the size of the coordinate ring \(\mathbb{K}[V]\) can also be measured using Hilbert series. First, by homogenizing the defining equations by adding one variable \(\gamma \in \mathbb{K}\), we embed the affine variety \(V \subset \mathbb{K}^d\) into the projective variety \(\mathcal{P} \subset \mathbb{K}^{d+1}\). The projective variety \(\mathcal{P}(h_1, \ldots, h_K) \subset \mathbb{K}^{d+1}\), defined by the homogeneous polynomials \(h_1, \ldots, h_K \in \mathbb{K}[(\bm{w}, \gamma)]\), is given by:
\begin{gather*}
\mathcal{P}(h_1, \ldots, h_K) \coloneqq \left\{ (\bm{w}, \gamma) \in \mathbb{K}^{d+1} : h_k(\bm{w}, \gamma) = 0,\ \forall k = 1,\ \ldots,\ K \right\}.
\end{gather*}

The dimension of the variety is also increased by one, i.e., \(d_{\mathcal{P}} = d_V + 1\). The coordinate ring \(\mathbb{K}[\mathcal{P}] = \mathbb{K}[(\bm{w, \gamma})]/I(\mathcal{P})\) of the projective variety \(\mathcal{P}\) can be decomposed into subgroups (called the graded coordinate ring) as follows:
\begin{align*}
    \mathbb{K}[\mathcal{P}] = \bigoplus_{\rho \in \mathbb{N}} S_{\rho},\ S_0 = \mathbb{K},
\end{align*}
where \(S_{\rho}\) is the set of homogeneous polynomials of degree \(\rho\) modulo the ideal \(I(\mathcal{P})\).
As a metric for the size of the coordinate ring \(\mathbb{K}[\mathcal{P}]\), the Hilbert function \(\operatorname{H}(\rho)\) and Hilbert-Poincaré series \(\operatorname{HS}(t)\) are defined as follows:
\begin{align*}
    \operatorname{H}(\rho) = \dim S_{\rho},\ \operatorname{HS}(t) = \sum_{\rho \in \mathbb{N}} \operatorname{H}(\rho)t^\rho = \frac{\prod_{k=1}^{K} (1 - t^{\rho_{k}})}{(1-t)^{d+1}},
\end{align*}
where \(\dim\) denotes the Krull dimension and \(\rho_{1}, \ldots, \rho_{K}\) are the degrees of the homogeneous polynomials \(h_1, \ldots, h_K\).

The Hilbert function represents the dimension of a "subspace" of the decomposed coordinate ring, and the Hilbert series is the generating function of the sequence of the Hilbert function, which is also a rational function with a pole at \(t = 1\). These measures indicate the growth of the dimension of the homogeneous components of the algebra with respect to the degree. According to the dimension theorem, the Krull dimension of the projective variety \(\mathcal{P}\) matches the order of the Hilbert series at the pole \(t = 1\), which is one of the most important results in commutative algebra.

Therefore, the dimension of the affine variety is defined using the Hilbert series, as follows:
\hilbertSeries*
Given the Gröbner basis of the ideal \(I(\mathcal{P})\), the Hilbert series can be easily computed, leading to an efficient estimation of the dimension of the affine variety \(d_V\).

\subsection{Regularity of Affine Varieties}
\label{app:regularity_of_affinev}
We informally define the concept of a regular set for real affine varieties, which is used in~\cref{sec:main-theorem} (for a formal definition, see Definition 2.1 in~\citep{zhang2023covering}). 

A affine variety \(V \subseteq \mathbb{R}^d\) is a \((\beta, d_V)\)-regular set if:
\begin{enumerate}
    \item For almost all affine planes \(L\) with \(\operatorname{codim}(L) \leq d_V\) in \(\mathbb{R}^d\), \(V \cap L\) has at most \(\beta\) path-connected components.
    \item For almost all affine planes \(L\) with \(\operatorname{codim}(L) > d_V\) in \(\mathbb{R}^d\), \(V \cap L\) is empty.
\end{enumerate}
The notion \(\operatorname{codim}\) represents the \textit{codimension}. For an affine subspace $L \subseteq \mathbb{R}^d$, its codimension is defined by $\operatorname{codim}(L)=d-d_L$.
Simply put, codimension is how many dimensions you are ``missing'' when comparing a smaller space inside a bigger space.

A regular set restricts the complexity of a variety \(V\). Intuitively, the complexity of \(V\) can be measured by the number of connected components in its cross sections. For instance, a complex shape may have cross sections that split into multiple connected components. The larger the number of connected components \(\beta\), the more complex the topology of \(V\).
Moreover, the dimension at which we slice the variety is also important. If the slice (affine plane) is large enough in dimension, \ie, the codimension is small ($<d_V$), then any intersection of the slice with $V$ is limited to at most $\beta$ connected pieces. Otherwise, the slice typically does not intersect \(V\) at all. For example, consider the circle $V = \{(x,y)\in \mathbb{R}^2 : x^2 + y^2 - 1=0\}$ .
A line ($\operatorname{codim}(L)=1$) intersects the circle in at most two points. For a single point ($\operatorname{codim}(L)=2$), almost all points do not lie in the circle; that is, intersections with higher codimension affine subspaces are almost empty. This implies that the circle is a $(2,1)$-regular set. 

\subsection{Covering number of Affine Varieties}
\begin{lemma}[\citet{zhang2023covering}]
\label{lem:cov}
Let \( V \subset \mathbb{R}^d \) be a \((\beta, d_V)\)-regular set in the ball \( \mathbb{B}_2(R) \) with the radius \( R \). Then for all \( \varepsilon \in (0, \text{diam}(V)] \),
\begin{gather}
\label{eq:lem:cov}
\log \mathcal{N}(V, \varepsilon, \|\cdot \|_2) \leq d_V \log \left(\frac{2Rd_V d}{\varepsilon}\right) + \log 2\beta.
\end{gather}
\end{lemma}
This upper bound is obtained by slicing the affine variety \( V \) with subspaces \( \{ L_{s} \}_{s \in \mathbb{N}} \) within \( \mathbb{R}^d \) and covering \( V \) with balls centered at the intersections of \( L_s \) and \( V \), \ie, \( V \subset \bigcup_{s} \bigcup_{v \in V \cap L_s} \mathbb{B}_2(v; \varepsilon) \). The covering for the two-dimensional case is illustrated in~\cref{fig:cover-ofaffineV}. The first term, \( \left(2Rd_V d/\varepsilon\right)^{d_V} \), represents the number of subspaces \( L_s \) needed to cover the entire space. It is mainly determined by the intrinsic dimension \( d_V \) of the affine variety, although it is still influenced by the ambient dimension \( d \). The quantity \( \beta \) in the second term denotes the number of intersections between a single subspace \( L \) and the variety \( V \), and represents the covering number of \( V \cap L \). Topologically, it corresponds to the Betti numbers of the affine variety, which informally represent the number of holes in \(V\). The upper bound on the quantity \( \beta \) is given, for example, by the Petrovskii-Oleinik-Milnor inequality~\citep{Petrovskii1949,oleinik1951estimates,milnor1964betti}. Specifically, an affine variety \( V \cap \mathbb{B}_2(R) \) defined by polynomials \( \{ p_k \}_{k\in [K]} \) of maximum degree \( \rho \) and the \(\ell_2\)-ball is \((\rho(2\rho-1)^{d+1}, d_V)\)-regular. This intuitively suggests that as the maximum degree of polynomials increases, the topology of the affine variety becomes more complex.

\section{Detailed Background on Problem Formulation}
\label{app:detail_background}
This appendix expands the formulation introduced in the main text, highlighting why a \emph{hybrid} physics–data approach is required.

\subsection{Governing System}
Let $\Omega\subset\mathbb{R}^{d}$ be a bounded domain with boundary $\partial\Omega$.  
For a differential operator $\mathscr{D}\!:L^{2}(\Omega)\to L^{2}(\Omega)$, the \emph{true} state
$f^{\ast}\!:\Omega\to\mathbb{R}$ satisfies the boundary‑value problem
\begin{align}
  \mathscr{D}[f^{\ast}] &= v && \text{in }\Omega, \label{eq:pde}\\
  f^{\ast} &= g && \text{on }\partial\Omega, \label{eq:bc}
\end{align}
where $v$ and $g$ are smooth but may be only \emph{partially observed} or inferred indirectly.

\subsection{Available Information}
In practice one seldom knows $v$ and $g$ exactly; instead one has:
\begin{itemize}
\item \textbf{Noisy pointwise observations.}\;
      A dataset $\{(x_{i},y_{i})\}_{i=1}^{N_{u}}$ with
      \(
        y_{i}=f^{\ast}(x_{i})+\varepsilon_{i},\;
        \varepsilon_{i}\sim\mathcal{N}(0,\sigma^{2}),
      \)
      where $x_{i}\in\Omega\cup\partial\Omega$.
\item \textbf{Weak‑form physics information.}\;
      Linear functionals
      $
        l_{k}(u) \coloneqq \langle u,\psi_{k}\rangle_{\mu_{k}},\;k=1,\dots,N_{r},
      $
      with trial functions $\psi_{k}\in C_{c}^{\infty}(\Omega)$
      and measures $\mu_{k}$,
      together with the corresponding targets
      $l_{k}(v)=l_{k}\!\left(\mathscr{D}[f^{\ast}]\right)$.
\end{itemize}

\subsection{Hybrid Surrogate Model}
\label{app:surrogate}
A representative hybrid approach is the
\emph{Physics‑Informed Neural Network} (PINN)~\citep{raissi2019physics}.
Given a neural surrogate $f_{\bm w}\!:\Omega\to\mathbb{R}$, its parameters $\bm w$ are obtained by minimising
\begin{align}
\mathcal{L}(\bm w)
  =\underbrace{\frac{1}{N_{u}}\sum_{i=1}^{N_{u}}
        \left|f_{\bm w}(x_{i})-y_{i}\right|^{2}}_{\text{data fidelity}}
  +\lambda\,
    \underbrace{\frac{1}{N_{r}}\sum_{k=1}^{N_{r}}
        \left|l_{k}\!\left(\mathscr{D}[f_{\bm w}]\right)-l_{k}(v)\right|^{2}}%
        _{\text{physics residual (weak form)}},
  \label{eq:hybrid-loss}
\end{align}
with hyper‑parameter $\lambda>0$ balancing empirical fit and physical consistency.

\paragraph{Connection to standard collocation method}
In the standard collocation method, the unified residual forms reduce to pointwise strong-form residuals:
specifically, for each $k$, we set
\begin{align*}
\psi_{k}(x)=1,\quad \mu_{k}=\delta_{x_{k}},
\end{align*}
where $\delta_{x_{k}}$ is the Dirac measure at collocation point $x_{k}$.
In this case, the linear functional $l_{k}$ becomes
\begin{align*}
l_{k}(u)=\langle u,\psi_{k}\rangle_{\mu_{k}}=\int u(x)\,\mathrm{d}\delta_{x_{k}}=u(x_{k}),
\end{align*}
and thus the physics term penalizes the pointwise physics residuals:
\begin{align*}
\sum_{k=1}^{N_{r}}\left|\mathscr{D}[f_{\bm w}](x_{k})-v(x_{k})\right|^{2}.
\end{align*}

\paragraph{Limiting regimes.}
\begin{itemize}
\item \emph{Pure data fitting:} $\lambda=0$ reduces \cref{eq:hybrid-loss} to standard supervised learning on $\mathcal{D}_{u}$.
\item \emph{Fully physics‑informed:}  
      If $\{l_{k}\}$ is dense and $N_{r}\to\infty$, the residual term enforces
      \cref{eq:pde} everywhere. 
\item \emph{Truly hybrid:}
      Finite $N_{r}$ with incomplete $\{l_{k}\}$—typical in engineering—captures partial physics,
      while the data term compensates for the missing information.
\end{itemize}

\subsection{Error decomposition and analysis}
\label{app:analysis setting}
To directly measure how physics-based inductive bias improves the generalization ability of the machine learning models, we fix the physics information as a known immutable prior. Let $\mathcal{H}_{\text{base}}$ denote the unrestricted hypothesis class (e.g.\ linear models or neural networks). We consider two ways of incorporating the physics prior into the base hypothesis class $\mathcal{H}_{\text{base}}$: 
\begin{itemize}
    \item \textbf{Hard setting:} Enforce the differential equation residual exactly by:
    \begin{align}
        \hat{f}_{\text{hard}} \in \argmin_{f \in \mathcal{H}_{\text{base}}} \sum_{k=1}^{N_r} |l_k(\mathscr{D}[f]) - l_k(v)|^2 := \mathcal{H}_{\text{hard}}.
    \end{align}
    \item \textbf{Soft setting:} Allow a relaxed residual tolerance:
    \begin{align}
        \hat{f}_{\text{soft}}(\varepsilon) \in \left\{f \in \mathcal{H}_{\text{base}} : \sum_{k=1}^{N_r} |l_k(\mathscr{D}[f]) - l_k(v)|^2 \le \varepsilon \right\} := \mathcal{H}_{\text{soft}}.
    \end{align}
\end{itemize}

For $\mathcal{H}\in\{\mathcal{H}_{\text{hard}},\mathcal{H}_{\text{soft}}(\varepsilon)\}$, let $\hat{f}\in\mathcal{H}$ be the learned solution and $f_{\mathcal{H}}^{\ast}:=\arg\min_{f\in\mathcal{H}}\|f^{\ast}-f\|$ the best attainable approximation. By the triangle inequality, we have:
\begin{align}
    \|f^* - \hat{f}\| \le \|f^* - f_{\mathcal{H}}^*\| + \|f_{\mathcal{H}}^* - \hat{f}\|.
\end{align}

In this decomposition:
\begin{itemize}
    \item The term $\|f^* - f_{\mathcal{H}}^*\|$ represents the approximation error: the best achievable error within the hypothesis class $\mathcal{H}$.
    \item The term $\|f_{\mathcal{H}}^* - \hat{f}\|$ represents the estimation error: the deviation due to finite data. 
\end{itemize}

Our primary focus is to derive bounds on the estimation error $\|f_{\mathcal{H}}^* - \hat{f}\|$,
thereby quantifying how well the learned solution converges to the best physics-constrained approximation.

In particular, our analysis centers on the hard constraint setting, where 
$\mathcal{H} = \mathcal{H}_{\text{hard}}$ on the linear base hypothesis
\[
   \mathcal{H}_{\text{base}} = 
   \left\{ f_{\bm{w}} = \bm{w}^{\top}\bm{\phi} :
      \bm{w} \in \mathbb{R}^{d},\ \phi_{j} \in \mathcal{B}\right\},
\]
with $\mathcal{B}$ a chosen basis.  
Under the hard constraint, the admissible hypothesis class amounts to restricting
the parameter vector $\bm{w}$ to lie on an affine variety induced by the PDE
residuals. More explicitly,
\[
   \mathcal{H}_{\text{hard}}
   = \Bigl\{ f_{\bm{w}} :
       \bm{w} \in \mathcal{V}(\mathscr{D},\mathcal{B},\mathcal{T}),\;
       \phi_{j}\in\mathcal{B}\Bigr\},
\]
where $\mathcal{V}(\mathscr{D},\mathcal{B},\mathcal{T})$ denotes the set of
coefficient vectors $\bm{w}$ satisfying the algebraic constraints generated by
the operator $\mathscr{D}$ acting on the basis $\mathcal{B}$ and tested against
the functionals $\mathcal{T}=\{l_{k}\}_{k=1}^{N_{r}}$.  


\subsection{Extension: Incomplete Operators with Learnable Parameters}
\label{app:extension}
The above framework can be generalized to settings where the governing
differential operator itself is only partially known and contains
learnable parameters.  
Formally, suppose that instead of a fixed operator
$\mathscr{D}:L^{2}(\Omega)\to L^{2}(\Omega)$,
we consider a parametric operator
\begin{align*}
   \mathscr{D}_{\bm{c}} : L^{2}(\Omega)\to L^{2}(\Omega), \qquad 
   \bm{c}\in\mathbb{R}^{m},
\end{align*}
where $\bm{c}$ denotes a vector of unknown coefficients to be
simultaneously estimated from data.  
In this case, the admissible hypothesis space is naturally
\begin{align*}
   \mathcal{H}_{\text{aug}}
   := \left\{ (f,\bm{c}) : f\in\mathcal{H}_{\text{base}},\;
        l_k\!\left(\mathscr{D}_{\bm{c}}[f]\right)=l_k(v),\;
        k=1,\dots,N_{r}\right\},
\end{align*}
defined as an \emph{augmented constraint set} over the joint variable
$(f,\bm{c})$.

The error decomposition then applies in this extended space:
for $(f^{\ast},\bm{c}^{\ast})$ denoting the best attainable pair in
$\mathcal{H}_{\text{aug}}$, the learned solution
$(\hat f,\hat{\bm{c}})$ satisfies
\begin{align*}
   \|f^{\ast}-\hat f\|
   \;\le\;
   \underbrace{\|f^{\ast}-f^{\ast}_{\mathcal{H}_{\text{aug}}}\|}_{\text{approximation error}}
   +\underbrace{\|f^{\ast}_{\mathcal{H}_{\text{aug}}}-\hat f\|}_{\text{estimation error}},
\end{align*}
with both terms now understood relative to the augmented parameter space.

\paragraph{Typical cases.}

\begin{itemize}
    \item \textbf{Unknown diffusion coefficient.}\;
    Consider the diffusion equation
    \(
      \partial_{t} u - c \Delta u = 0,
    \)
    where the diffusion constant $c>0$ is unknown.
    Here $\bm{c}=(c)$ is a scalar parameter.  
    Under the hard constraint, the admissible hypothesis class can be written as
    \[
      \mathcal{H}_{\text{aug}}
      = \Bigl\{ (f_{\bm{w}}, c)  :
          (\bm{w},c) \in 
          \mathcal{V}\!\bigl(\partial_{t}f_{\bm{w}} - c\Delta f_{\bm{w}},
                              \mathcal{B},\mathcal{T}\bigr),
          \ \phi_{j}\in\mathcal{B}\Bigr\},
    \]
    where $\mathcal{V}(\cdot)$ denotes the algebraic variety of
    coefficient–parameter pairs $(\bm{w},c)$ that satisfy the residual
    constraints induced by $\mathcal{T}=\{l_{k}\}_{k=1}^{N_{r}}$.
    Thus both approximation and estimation errors are quantified in this
    augmented parameter space.

    \item \textbf{Unknown diffusion term (learned surrogate).}\;
    In cases where the diffusion operator itself is not specified, one
    may introduce a surrogate $v_{\theta}$ to represent its action.
    The PDE constraint becomes
    \[
      \partial_{t} f_{\bm{w}} - v_{\theta} = 0,
    \]
    leading to the hypothesis class
    \[
      \mathcal{H}_{\text{aug}}
      = \Bigl\{ (f_{\bm{w}}, v_{\theta}) :
          (\bm{w},\theta) \in 
          \mathcal{V}\!\bigl(\partial_{t}f_{\bm{w}} - v_{\theta},
                              \mathcal{B},\mathcal{T}\bigr),
          \ \phi_{j}\in\mathcal{B}\Bigr\}.
    \]
    Here $v_{\theta}$ serves as a learnable proxy for the unknown
    diffusion term.  
    Our error decomposition applies verbatim in this augmented parameter
    space, with approximation error defined relative to the best attainable
    pair $(\bm{w}^{\ast},\theta^{\ast})$ and estimation error measuring the
    deviation of the learned $(\hat{\bm{w}},\hat{\theta})$ from this target.
\end{itemize}

In summary, by enlarging the hypothesis space to include both explicit
unknown coefficients and implicit unknown operator surrogates, the
proposed error decomposition continues to hold, thereby providing a
principled means of quantifying generalization in operator-learning
settings with incomplete physics.

% \subsection{Relation to Prior Theoretical Work on Physics-Informed Models}

\section{Proof for~\cref{thm:main}}
\label{app:proof_thm:main}
\main*
\begin{proof}[Proof]
\textbf{Step 1:} We first upper bound the prediction error by a term that represents the supremum of a empirical process in the metric space of the affine variety.
Using~\cref{lem:bound_Phi_beta}, we get:
\begin{align*}
\|\bm{\Phi}(\bm{w}^* - \hat{\bm{w}})\|_2^2 \leq 2 \bm{\varepsilon}^{\top} \bm{\Phi}(\bm{w}^* - \hat{\bm{w}}).
\end{align*}
We denote \(\rx_{\bm{w}} \coloneqq \bm{\varepsilon}^{\top} \bm{\Phi}(\bm{w}-\hat{\bm{w}})\) as the random process in the metric space \((\mathcal{V}_R , \|\cdot\|_2)\). Note that the estimator \(\hat{\bm{w}}\) is a random variable depending on the parameter \(\bm{w}\) and the noise \(\bm{\varepsilon}\). Then, the minimax risk is bounded as follows.
\begin{align}
    \label{eq:bound_minimax_risk}
    \min_{\hat{\bm{w}}} \max_{\bm{w}^* \in \mathcal{V}_R} \|\hat{\bm{w}} - \bm{w}^*\|_2^2 \leq \min_{\hat{\bm{w}}} \max_{\bm{w}^* \in \mathcal{V}_R} \frac{\eta^{-1}}{n} \|\bm{\Phi} (\hat{\bm{w}} - \bm{w}^*)\|_2^2 \leq \frac{2}{n} \eta^{-1} \sup_{\bm{w} \in \mathcal{V}_R} \rx_{\bm{w}}.
\end{align}
The first inequality holds by~\cref{asm:lowereig}. 

\textbf{Step 2:} Next, we calculate the supremum of the empirical process \(\rx_{\bm{w}}\) using the covering number. 
For all \(\bm{w}_1, \bm{w}_2 \in \mathcal{V}_R\), it is shown that the variable \(\rx_{\bm{w}_1} - \rx_{\bm{w}_2}\) has sub-Gaussian increments with respect to the metric \(\|\cdot\|_2\):
\begin{align}
    \label{eq:Gaussian_increments}
    \begin{aligned}
    \rx_{\bm{w}_1}- \rx_{\bm{w}_2} &= \sum_{i=1}^n \varepsilon_i ((\bm{w}_1-\hat{\bm{w}}_1)-(\bm{w}_2-\hat{\bm{w}}_2))^{\top}\bm{\phi}(x_i) \\
    &\leq \sum_{i=1}^n \varepsilon_i \|(\bm{w}_1-\bm{w}_2)-(\hat{\bm{w}_1}-\hat{\bm{w}}_2)\|_2\|\bm{\phi}(x_i)\|_2 \\
    &\leq \sum_{i=1}^n \varepsilon_i \left(\|\bm{w}_1-\bm{w}_2\|_2 +\|\hat{\bm{w}_1}-\hat{\bm{w}}_2\|_2\right) M \\
    &\leq \Gamma \|\bm{w}_1-\bm{w}_2\|_2 M \re,
    \end{aligned}
\end{align}
where \(\re\) is the zero-mean Gaussian random variable with variance \(n\sigma^2\). The second inequality holds by the Cauchy-Schwarz inequality and the third holds by the triangle inequality and~\cref{asm:bounded_phi}. The last inequality holds by~\cref{asm:stability_est}.

From~\cref{eq:Gaussian_increments}, the random process \(\rx_{\bm{w}_1} - \rx_{\bm{w}_2}\) has sub-Gaussian increments as follows.
\begin{align*}
    \|\rx_{\bm{w}_1}-\rx_{\bm{w}_2}\|_{\psi_2} \leq \sqrt{n}\sigma M \Gamma \|Z\|_{\psi_2} \| \bm{w}_1- \bm{w}_2\|_2,
\end{align*}
where \(Z\) is the standard Gaussian random variable and \(\|\cdot\|_{\psi_2}\) is the sub-Gaussian norm.
For the centered random process \(\rz_{\bm{w}} \coloneqq \rx_{\bm{w}} - \mathbb{E}[\rx_{\bm{w}}] \), \(\|\rz_{\bm{w}_1}-\rz_{\bm{w}_2}\|_{\psi_2} \lesssim \|\rx_{\bm{w}_1}-\rx_{\bm{w}_2}\|_{\psi_2}\) holds because \(\|\rx_{\bm{w}_1}-\rx_{\bm{w}_2}\|_{\psi_2}\) is sub-Gaussian.

Using~\cref{lem:sup_z}, we obtain the following bound with some constant $C_0$:
\begin{align}
    \label{eq:sup_z}
    \mathbb{E}\sup_{\bm{w} \in \mathcal{V}_R} \rz_{\bm{w}}\leq C_0 \sqrt{n}\sigma M \Gamma R \left( \sqrt{d_{\mathcal{V}}\log d_{\mathcal{V}}d} + \sqrt{\log 2\beta} \right).
\end{align}
Next, using Dudley's integral tail bound, we have:
\begin{align*}
    \Pr\left(\sup_{\bm{w} \in \mathcal{V}_R} \rz_{\bm{w}} \leq  \mathbb{E}\sup_{\bm{w} \in \mathcal{V}_R} \rz_{\bm{w}} + C_0 \sqrt{n}\sigma M \Gamma 2R \sqrt{\log (\delta/2)}\right) \geq 1 - \delta.
\end{align*}
By incorporating the non-centered process \(\rx_{\bm{w}}\), we obtain:
\begin{align}
    \label{eq:tail_xb}
    \Pr\left(\sup_{\bm{w} \in \mathcal{V}_R} \rx_{\bm{w}} \leq \sup_{\bm{w} \in \mathcal{V}_R} \left|\mathbb{E}[\rx_{\bm{w}}]\right| + \mathbb{E}\sup_{\bm{w} \in \mathcal{V}_R} \rz_{\bm{w}} + C_0 \sqrt{n}\sigma M \Gamma 2R \sqrt{\log (\delta/2)}\right) \geq 1 - \delta.
\end{align}
To bound \(\mathbb{E}[\rx_{\bm{w}}]\), we note that:
\begin{align}
    \begin{aligned}
    \label{eq:e_xb}
    \mathbb{E}[\rx_{\bm{w}}] &= \mathbb{E}\left[ \bm{\varepsilon}^{\top}  \bm{\Phi} \left(\bm{w}- \hat{\bm{w}}\right) \right] \\
    &= \mathbb{E}\left[ \bm{\varepsilon}^{\top} \bm{\Phi} \hat{\bm{w}} \right] \\
    &\leq \sqrt{\mathbb{E}\left[ \|\bm{\varepsilon}^{\top} \bm{\Phi}\|_2^2\right]}\sqrt{\mathbb{E}\left[\left\| \hat{\bm{w}} \right\|^2_2\right]} \\
    &\leq \sigma \sqrt{\sum_{j=1}^d \sum_{i=1}^n |\phi_j(x_i)|^2} R \\
    &= \sigma \sqrt{n} M R
    \end{aligned}
\end{align}
Here, the third inequality follows from the Cauchy-Schwarz inequality, and the fourth inequality is derived from the fact that \(|\bm{\varepsilon}^{\top} \bm{\Phi}_j|^2 / (\sigma \|\bm{\Phi}_j\|_2)^2\) follows a chi-squared distribution with \(1\) degrees of freedom and \(\hat{\bm{w}} \in \mathcal{V}_R\). 

By combining \cref{eq:sup_z}, \cref{eq:tail_xb}, and \cref{eq:e_xb}, we obtain the following bound with some constant \(C\):
\begin{align*}
    \Pr\left(\frac{2}{n} \sup_{\bm{w} \in \mathcal{V}_R} \rx_{\bm{w}} \leq C \sigma M \Gamma R
    \left( \sqrt{\frac{d_{\mathcal{V}}\log d_{\mathcal{V}}d}{n}} + \sqrt{\frac{\log 2\beta}{n}} + 2\sqrt{\frac{\log (\delta/2)}{n}}\right) \right) \geq 1 - \delta.
\end{align*}
This completes the proof.
\end{proof}

\begin{lemma}
\label{lem:bound_Phi_beta}
Let \(\hat{\bm{w}}\) be a minimizer of the following optimization problem:
\begin{gather}
\label{eq:lin-problem-ball-app}
\hat{\bm{w}} = \argmin_{\bm{w} \in \mathcal{V}_R} \frac{1}{n} \| \bm{y} - \bm{\Phi} \bm{w} \|^2_2,
\end{gather}
where \( \mathcal{V}_R = \mathcal{V}(\mathscr{D}, \mathcal{B}, \mathcal{T}) \cap \mathbb{B}_2(R) \) is the affine variety constrained with the \(\ell_2\)-ball, \(\bm{y} = \bm{\Phi}\bm{w}^* + \bm{\varepsilon}\) is the observed vector, \(\bm{\Phi}\) is the design matrix, \(\bm{w}^* \in \mathcal{V}_R\) is the true parameter vector, and \(\bm{\varepsilon} = [\varepsilon_1, \ldots, \varepsilon_n]^{\top}\) is the noise vector with each \(\varepsilon_i\) independently following a zero-mean Gaussian distribution. Then, under these conditions, we have:
\begin{align}
\label{eq:lemma_basic_bound}
\|\bm{\Phi}(\bm{w}^* - \hat{\bm{w}})\|_2^2 \leq 2 \bm{\varepsilon}^{\top} \bm{\Phi}(\bm{w}^* - \hat{\bm{w}}).
\end{align}
\begin{proof}[Proof]
Since \(\hat{\bm{w}}\) is a minimizer of~\cref{eq:lin-problem-ball-app}, we have:
\begin{align*}
    \|\bm{y} - \bm{\Phi}\hat{\bm{w}}\|_2^2 \leq \|\bm{y} - \bm{\Phi}\bm{w}^*\|_2^2 = \|\bm{\varepsilon}\|_2^2.
\end{align*}
The left-hand side can be expanded as:
\begin{align*}
    \|\bm{y} - \bm{\Phi}\hat{\bm{w}}\|_2^2 &= \|\bm{y} - \bm{\Phi}\bm{w}^* + \bm{\Phi}\bm{w}^* - \bm{\Phi}\hat{\bm{w}}\|_2^2 \\
    &= \|\bm{\varepsilon} - \bm{\Phi}(\bm{w}^* - \hat{\bm{w}})\|_2^2.
\end{align*}
Thus, we have:
\begin{align*}
    \|\bm{\varepsilon} - \bm{\Phi}(\bm{w}^* - \hat{\bm{w}})\|_2^2 \leq \|\bm{\varepsilon}\|_2^2.
\end{align*}
Expanding the left-hand side, we get:
\begin{align*}
    \|\bm{\varepsilon} - \bm{\Phi}(\bm{w}^* - \hat{\bm{w}})\|_2^2 &= \|\bm{\varepsilon}\|_2^2 - 2 \bm{\varepsilon}^{\top} \bm{\Phi}(\bm{w}^* - \hat{\bm{w}}) + \|\bm{\Phi}(\bm{w}^* - \hat{\bm{w}})\|_2^2.
\end{align*}
Subtracting \(\|\bm{\varepsilon}\|_2^2\) from both sides, we obtain:
\begin{align*}
    \|\bm{\Phi}(\bm{w}^* - \hat{\bm{w}})\|_2^2 \leq 2 \bm{\varepsilon}^{\top} \bm{\Phi}(\bm{w}^* - \hat{\bm{w}}).
\end{align*}
This completes the proof.
\end{proof}
\end{lemma}

\begin{lemma}
\label{lem:sup_z}
Let \(\rz_{\bm{w}}\) be the zero-mean random process in the metric space \((\mathcal{V}_R, \|\cdot\|_2)\), which have the following sub-Gaussian increments. For all \(\bm{w}_1, \bm{w}_2 \in \mathcal{V}_R\),
\begin{align*}
    \| \rz_{\bm{w}_1} - \rz_{\bm{w}_2} \|_{\psi_2} \leq A \|\bm{w}_1 - \bm{w}_2\|_2,
\end{align*}
where \(\|\cdot\|_{\psi_2}\) is the sub-Gaussian norm, \(A\) is a positive constant. Then, the expectation of the supremum of the process can be bounded as follows.
\begin{align*}
    \mathbb{E}\sup_{\bm{w} \in \mathcal{V}_R} \rz_{\bm{w}}\leq C A R \left( \sqrt{d_{\mathcal{V}}\log d_{\mathcal{V}}d} + \sqrt{\log 2\beta} \right),
\end{align*}
where \(C\) is positive constant. 

\begin{proof}
Using Dudley’s integral inequality~\citep{dudley1967sizes} to the zero-mean random process:
\begin{gather}
    \label{eq:cov_bound_for_GP}
    \mathbb{E} \sup_{\bm{w} \in \mathcal{V}_R} \rz_{\bm{w}} \leq C_0 A \int_0^{\infty} \sqrt{\log \mathcal{N}(\mathcal{V}_R, \varepsilon, \|\cdot\|_2)} \mathrm{d}\varepsilon.
\end{gather}

Since the set \(\mathcal{V}_R\) is \((\beta, d_{\mathcal{V}}
)\) regular set from Lemma 2.13 by~\citet{zhang2023covering}, \Cref{lem:cov} shows the upper bound of the covering number for any $\varepsilon \in (0, 2R]$ as follows. 
\begin{align*}
    \log \mathcal{N}(\mathcal{V}_R, \varepsilon, \|\cdot\|_2) \leq d_{\mathcal{V}} \log \left(\frac{2Rd_{\mathcal{V}} d}{\varepsilon}\right) + \log 2\beta.
\end{align*}

We substitute the above inequality to~\cref{eq:cov_bound_for_GP}:
\begin{align*}
    \mathbb{E}\sup_{\bm{w} \in \mathcal{V}_R} \rz_{\bm{w}}\leq  C_0A\left( \sqrt{d_{\mathcal{V}}}\int^{\infty}_0 \sqrt{\log \left(\frac{2Rd_{\mathcal{V}}d}{\varepsilon}\right)} \mathrm{d}\varepsilon + 2R \sqrt{\log 2\beta} \right).
\end{align*}
The integral in the first term can be calculated using substitution and integration by parts. Let
\begin{align*}
I \coloneqq \int^{\infty}_0 \sqrt{\log \left(\frac{2Rd_{\mathcal{V}}d}{\varepsilon}\right)} \mathrm{d}\varepsilon = \int^{2R}_0 \sqrt{\log \left(\frac{2Rd_{\mathcal{V}}d}{\varepsilon}\right)} \mathrm{d}\varepsilon.
\end{align*}
We substitute \(\chi \coloneqq 2Rd_{\mathcal{V}}d,\ u := \log (\chi/\varepsilon)\) into the integral:
\begin{align*}
    I = \int_{\infty}^{\log d_{\mathcal{V}}d} u^{1/2}(-\chi e^{-u})\mathrm{d}u.
\end{align*}
To solve the above integral, we use the formula for integration by parts:
\begin{align*}
    \begin{aligned}
    I &= -\chi \left([-u^{1/2}e^{-u}]^{\log d_{\mathcal{V}}d}_{\infty} + \frac{1}{2} \int^{\log d_{\mathcal{V}}d}_{\infty} u^{-1/2}e^{-u}\mathrm{d}u \right) \\
    &= 2R \sqrt{\log d_{\mathcal{V}}d} + Rd_{\mathcal{V}}d \int^{\infty}_{\log d_{\mathcal{V}}d} u^{-1/2}e^{-u}\mathrm{d}u.
    \end{aligned}
\end{align*}
The integral in the second term can be upper bounded as follows.
\begin{align*}
    \int^{\infty}_{\log d_{\mathcal{V}}d} u^{-1/2}e^{-u}\mathrm{d}u \leq \int^{\infty}_{\log d_{\mathcal{V}}d} e^{-u}\mathrm{d}u = [-e^{-u}]^{\infty}_{\log d_{\mathcal{V}}d} = (d_{\mathcal{V}}d)^{-1}.
\end{align*}

We obtain the following bound with some constant \(C\). 
\begin{align*}
    \mathbb{E}\sup_{\bm{w} \in \mathcal{V}_R} \rz_{\bm{w}}\leq C A R \left( \sqrt{d_{\mathcal{V}}\log d_{\mathcal{V}}d} + \sqrt{\log 2\beta} \right).
\end{align*}
\end{proof}
\end{lemma}

\section{Proof for~\cref{thm:pi-kernel} and~\cref{thm:lb-dim}}
\label{app:proof_thm:pi-kernel}
\eff*
\begin{proof}[Proof]
    From Theorem 4.2 in~\citep{doumeche2024physics} and Equation 15 in~\citep{doumeche2024physics2}, the effective dimension is bounded as follows:
    \begin{align}
        \label{eq:eff-bound-ref}
        d_{\mathrm{eff}}(\xi, \nu) \lesssim \sum_{\alpha \in \sigma(\mB\mM^{-1}\mB)} \frac{1}{1 + \alpha^{-1}} \leq \sum_{\alpha \in \sigma(\mM^{-1})} \frac{1}{1 + \alpha^{-1}},
    \end{align}
    where \(\mM \coloneqq \xi \mI + \nu \mD^{\top}\mT\mD \in \mathbb{R}^{d \times d}\) and \(\mB \in \mathbb{R}^{d \times d}\) is the Gram matrix of the basis functions, i.e., \(\emB_{j, j'} = \langle \phi_j, \phi_{j'} \rangle_{\mu} \) for all \(\phi_j, \phi_{j'} \in \mathcal{B}\).
    
    Since the matrix \(\mD^{\top}\mT\mD\) is positive semi-definite, the eigenvalues of the matrix \(\mM\) in ascending order \(\sigma_j(\cdot)\) are given by
    \begin{align*}
        \sigma_j(\mM) = \left\{\begin{array}{cc}
            \xi & (j = 1, \ldots, d_{\mathcal{V}}) \\
            \xi + \nu \alpha_j & (d_{\mathcal{V}} < j)
        \end{array} \right..
    \end{align*}
    Therefore, the matrix \(\mM\) is positive definite, and the eigenvalues of \(\mM^{-1}\) are \( \alpha^{-1} \) for all \(\alpha \in \sigma(\mM)\). Combining this with \cref{eq:eff-bound-ref}, we obtain the first inequality. The second inequality is obtained when \(\nu = 0\).
\end{proof}

\lowerbound*
\begin{proof}
The point \( \bm{w} =\bm{0} \) lies on \( \mathcal{V}(\mathscr{D}) \), and if \( \mathscr{L} \neq 0 \), it is not singular. The Jacobian rank of polynomials \( p_k(\bm{w}) = \langle \mathscr{D}[\bm{w}^{\top} \bm{\phi}], \psi_k \rangle_{\mu_k} \) in \( \bm{w} = \bm{0} \) is equal to \(d - d_{\mathcal{V}(\mathscr{L})}\). By \cref{def:tangent-space}, we have \(d_{\mathcal{V}(\mathscr{L})} \leq  d_{\mathcal{V}(\mathscr{D})}\).
\end{proof}


\section{Generalization Bounds via Rademacher Complexity}
\label{app:Rademacher}
In this section, we derive a generalization bound for our regression problem by controlling the Rademacher complexity of the hypothesis class under additional mild assumptions.

First, we introduce the hypothesis space constrained by an $\ell_{p}$–ball of radius $R$.

\begin{definition}[Hypothesis Space]
  Let $\mathcal{B}$ be a fixed set of basis functions and
  $$\mathbb{B}_{p}(R) = \left\{\,\bm{w}\in\mathbb{R}^{d}\;\bigm|\;\|\bm{w}\|_{p}\le R\right\}.$$ 
  We define the hypothesis space
  \begin{align*}
    \mathcal{H}
    &= \left\{\,f_{\bm{w}}(x) = \bm{w}^{\top}\boldsymbol{\phi}(x)
       \;\bigm|\;\bm{w}\in\mathcal{V}\cap\mathbb{B}_{p}(R),\;
           \phi_{j}\in\mathcal{B}\right\},
  \end{align*}
  where $\mathcal{V}(\mathscr{D},\mathcal{B},\mathcal{T})$ is the affine variety defined by the differential‐operator constraints.
\end{definition}

Next, we recall the definition of Rademacher complexity.

\begin{definition}[Rademacher Complexity]
  Given a function class $\mathcal{F}$ and a sample
  $S = \{x_{1},\dots,x_{n}\}$ drawn i.i.d., the (empirical) Rademacher complexity is
  \begin{align*}
    \widehat{\mathfrak{R}}_{S}(\mathcal{F})
    &= \frac{1}{n}\,\mathbb{E}_{\boldsymbol{\tau}}\Biggl[\,
       \sup_{f\in\mathcal{F}}
       \left|\sum_{i=1}^{n}\tau_{i}\,f(x_{i})\right|\Biggr],
  \end{align*}
  where $\tau_{1},\dots,\tau_{n}$ are independent Rademacher variables taking values in $\{\pm1\}$.  The Rademacher complexity is its expectation over $S$.
\end{definition}

To proceed, we impose the following boundedness assumptions.

\begin{assumption}[Bounded Target]\label{assump:bounded_target}
  The true regression function $f^{*}:\mathbb{R}^{m}\to\mathbb{R}$ satisfies
  $$\|f^{*}\|_{\infty}\le F_{\max}\,. $$
\end{assumption}

\begin{assumption}[Bounded Hypotheses and Features]\label{assump:bounded_H}
  Let $p,q\in[1,\infty]$ with $1/p+1/q=1$.  Every hypothesis $f_{\bm{w}}\in\mathcal{H}$ satisfies
  \begin{align*}
  \|f_{\bm{w}}\|_{\infty}\le F_{\max},
  \end{align*}
  and each feature vector $\boldsymbol{\phi}(x)$ satisfies
  \begin{align*}
  \|\boldsymbol{\phi}(x)\|_{q}\le M\quad\forall\,x\,. 
  \end{align*}
\end{assumption}

We now establish a key concentration lemma for the squared‐loss deviations.
\begin{lemma}
\label{lem:bounded_g}
Let 
\(g_{\bm w}(x,y) := (f_{\bm w}(x)-y)^2 \;-\; \mathbb{E}_{(X,Y)}\left[(f_{\bm w}(X)-Y)^2\right]\) for $f_{\bm w} \in \mathcal{H}$. Under \cref{assump:bounded_target} and \cref{assump:bounded_H} and 
assume the noise \(\epsilon = Y - f^*(X)\) is sub‐Gaussian with proxy variance \(\sigma^2\). Then there exists a constant \(C >0\) (independent of \(\bm w\)) such that
\begin{align*}
\|g_{\bm w}\|_{\psi_{1}}
\le 2(4F_{\max}+\sigma)^{2},
\qquad
\sup_{\bm w}\mathbb{E}\left[g_{\bm w}(X,Y)^{2}\right]
\le C(4F_{\max}+\sigma)^4
\end{align*}
\end{lemma}

\begin{proof}
By the triangle inequality for the sub‐Gaussian norm,
\begin{align*}
\|f_{\bm w}(X) - Y\|_{\psi_{2}}
&\le \|f_{\bm w}(X)-f^*(X)\|_{\psi_{2}}
   + \|\epsilon\|_{\psi_{2}}\\
&\le 4F_{\max} + \sigma
\end{align*}
Using the identity \(\|Z^2\|_{\psi_{1}} = \|Z\|_{\psi_{2}}^2\) for any $Z$, we get
\begin{align*}
\|(f_{\bm w}(X)-Y)^2\|_{\psi_{1}} \le (4F_{\max} + \sigma)^2.
\end{align*}
By the triangle inequality in the \(\psi_1\)–norm,
\begin{align*}
\|g_{\bm w}(X,Y)\|_{\psi_{1}}
&\le \|(f_{\bm w}(X)-Y)^2\|_{\psi_{1}}
     + \|\mathbb{E}(f_{\bm w}(X)-Y)^2\|_{\psi_{1}}\\
&\le 2(4F_{\max} + \sigma)^2.
\end{align*}

Finally, to bound the second moment,
\begin{align*}
\mathbb{E}\left[g_{\bm w}(X,Y)^2\right]= \mathrm{Var}\left(g_{\bm w}(X,Y)\right) \leq C(4F_{\max}+\sigma)^4
\end{align*}
for suitable constant \(C>0\). This completes the proof.
\end{proof}

\begin{lemma}[Maximum of sub-exponential random variables]\label{lem:max_subexp}
Let \(X_1,\dots,X_n\) be independent, identically distributed sub-exponential
random variables satisfying \(\|X_i\|_{\psi_1}\le\nu\) for every
\(i \in \{1,\dots,n\}\).  
Then there exists an absolute constant \(C>0\) such that
\begin{align*}
    \left\|\max_{1\le i\le n}\lvert X_i\rvert\right\|_{\psi_1}
    \;\le\; C\,\nu\log n .
\end{align*}
\end{lemma}

\begin{proof}
From the assumption \( \|X_i\|_{\psi_1} \leq \nu \), we have the standard sub-exponential tail bound:
\begin{align*}
    \Pr\left\{ \max_{1 \le i \le n} X_i \ge t \right\}
    &\le
    2\exp\left(
        -\frac{1}{2} \min\left\{
            \frac{t^2}{\nu^2},\;
            \frac{t}{\nu}
        \right\}
    \right).
\end{align*}

Moreover, for all \( t \ge 0 \), we can uniformly bound the maximum via a union bound:
\begin{align}
    \Pr\left\{ \max_{1 \le i \le n} X_i \ge t \right\}
    &= \Pr\left[ \bigcup_{i=1}^{n} \{ X_i \ge t \} \right]
    \le \sum_{i=1}^{n} \Pr\{ X_i \ge t \} \notag \\
    &\le
    2n \exp\left(
        -\min\left\{
            \frac{t^2}{2\nu^{2}},\;
            \frac{t}{2\nu}
        \right\}
    \right).
    \label{eq:union_bound}
\end{align}

Now consider two cases based on the relation between \( n \) and the exponent.

\smallskip
\textbf{Case 1:} Suppose
\begin{align*}
    n \le \exp\left(
        \frac{1}{2} \min\left\{
            \frac{t^2}{2\nu^{2}},\;
            \frac{t}{2\nu}
        \right\}
    \right).
\end{align*}
Then, by inequality \eqref{eq:union_bound}, we have
\begin{align*}
    \Pr\left\{ \max_{i} X_i \ge t \right\}
    \le
    2\exp\left(
        -\frac{1}{2} \min\left\{
            \frac{t^2}{2\nu^2},\;
            \frac{t}{2\nu}
        \right\}
    \right)
    \le
    2\exp\left(
        -\frac{1}{3\log n}
        \min\left\{
            \frac{t^2}{2\nu^2},\;
            \frac{t}{2\nu}
        \right\}
    \right),
\end{align*}
which already provides a stronger tail decay than we ultimately require.

\smallskip
\textbf{Case 2:} Suppose
\begin{align*}
    n > \exp\left(
        \frac{1}{2} \min\left\{
            \frac{t^2}{2\nu^{2}},\;
            \frac{t}{2\nu}
        \right\}
    \right).
\end{align*}
Then, the probability on the right-hand side of \cref{eq:union_bound} satisfies
\begin{align*}
    \Pr\left\{ \max_{1 \leq i \leq n} X_i \ge t \right\}
    \le
    2\exp\left(
        -\frac{1}{3\log n}
        \min\left\{
            \frac{t^2}{2\nu^2},\;
            \frac{t}{2\nu}
        \right\}
    \right)
    > 2e^{-2/3} > 1,
\end{align*}
which is vacuously bounded above by 1. Thus, it does not affect the validity of our bound.

In either case, we obtain the uniform upper bound
\begin{align*}
    \Pr\left\{ \max_{1 \leq i \leq n} X_i \ge t \right\}
    &\le
    2\exp\left(
        -\frac{1}{3\log n}
        \min\left\{
            \frac{t^2}{2\nu^2},\;
            \frac{t}{2\nu}
        \right\}
    \right) \\
    &=
    2\exp\left(
        -\frac{1}{3}
        \min\left\{
            \frac{t^2}{2(\nu\sqrt{\log n})^2},\;
            \frac{t}{2(\nu\log n)}
        \right\}
    \right).
\end{align*}

This tail bound implies that
\begin{align*}
    \left\| \max_{1 \leq i \leq n} X_i \right\|_{\psi_1} \leq C \nu \log n,
\end{align*}
for some universal constant \( C > 0 \).
\end{proof}
\begin{lemma}
\label{lem:contraction_rad}
Let $\ell(u,y) = (u - y)^2$.  The Rademacher complexity of the composite class
$\ell \circ \mathcal{H} = \{(f_{\bm w}(x) - y)^2 : f_{\bm w}\in\mathcal{H}\}$ satisfies
\[
\mathfrak{R}(\ell\circ\mathcal{H})
\;\le\;
\left(4F_{\max} + 2\sigma\sqrt{2/\pi}\right)\,
\mathfrak{R}(\mathcal{H}),
\]
where $F_{\max}$ is a uniform bound on $|f_{\bm w}(x)|$ and $\sigma^2$ is the noise variance. 
\end{lemma}

\begin{proof}
Define for each example $(x_i,y_i)$ the function
\begin{align*}
\phi_i(u) = (u - y_i)^2.
\end{align*}
For any $u,v\in\mathbb{R}$,
\begin{align*}
\left|\phi_i(u) - \phi_i(v)\right|
&= \left|u^2 - v^2 + 2y_i\,(v - u)\right|\nonumber\\
&= \left|(u - v)\,(u + v - 2y_i)\right|\nonumber\\
&\le \left(|u| + |v| + 2|y_i|\right)\,\left|u - v\right|.
\end{align*}
Since $|f_{\bm w}(x)|\le F_{\max}$ for all $x$ and $|y_i|\le F_{\max} + |\varepsilon_i|$, it follows that
\begin{align*}
\left|\phi_i(u) - \phi_i(v)\right|
\;\le\;
\left(4F_{\max} + 2|\varepsilon_i|\right)\,\left|u - v\right|.
\end{align*}
Applying the Rademacher contraction lemma to the empirical Rademacher complexity $\widehat{\mathfrak{R}}_{S}$,
\begin{align*}
\mathfrak{R}(\ell\circ\mathcal{H})
&= \mathbb{E}_{X,\varepsilon_{1:n}}\left[\widehat{\mathfrak{R}}_{S}(\ell\circ\mathcal{H})\right]\nonumber\\
&\le \mathbb{E}_{X,\varepsilon}\left[(4F_{\max} + 2|\varepsilon|)\,\widehat{\mathfrak{R}}_{S}(\mathcal{H})\right]\nonumber\\
&= \left(4F_{\max} + 2\,\mathbb{E}[|\varepsilon|]\right)\,
   \mathbb{E}_{X}\left[\widehat{\mathfrak{R}}_{S}(\mathcal{H})\right].
\end{align*}
Since $\mathbb{E}_{\varepsilon}[|\varepsilon|]\le\sigma\sqrt{2/\pi}$ for Gaussian noise,
\begin{align*}
\mathfrak{R}(\ell\circ\mathcal{H})
\;\le\;
\left(4F_{\max} + 2\sigma\sqrt{2/\pi}\right)\,
\mathfrak{R}(\mathcal{H}).
\end{align*}
This completes the proof.  
\end{proof}

\begin{lemma}[Dudley Integral Bound]
\label{lem:dudley_rad}
Let $\mathcal{H}$ be as above, with $d_{\mathcal{V}} = \dim(\mathcal{V}_{p,R})$.  Then there exists a constant $C>0$ such that
\begin{align*}
\mathfrak{R}(\mathcal{H})
&\le
C\,F_{\max}\,M
\left(
\sqrt{\,\frac{d_{\mathcal{V}}}{n}\;\ln\left(\frac{2R\,d_{\mathcal{V}}\,d}{M\,F_{\max}}\right)}
+\sqrt{\frac{\ln(2\beta)}{n}}
\right).
\end{align*}
\end{lemma}

\begin{proof}
First, for any $\bm w_1,\bm w_2\in\mathcal{V}_{p,R}$ and corresponding $f_{\bm w_1},f_{\bm w_2}\in\mathcal{H}$, by Hölder’s inequality and $\|\bm\phi\|_{q}\le M$, we have for all $\varepsilon>0$,
\begin{align*}
\|\bm w_1-\bm w_2\|_{p}\le\varepsilon
&\;\Longrightarrow\;
\|f_{\bm w_1}-f_{\bm w_2}\|_{1}
=\|\bm w_1^{\top}\bm\phi - \bm w_2^{\top}\bm\phi\|_{1}\\
&\le
\|\bm w_1-\bm w_2\|_{p}\,\|\bm\phi\|_{q}
\le\varepsilon\,M.
\end{align*}
Hence
\begin{align*}
\mathcal{N}\left(\mathcal{H},\varepsilon M,\|\cdot\|_{1}\right)
&\le
\mathcal{N}\left(\mathcal{V}_{p,R},\varepsilon,\|\cdot\|_{p}\right).
\end{align*}

By Dudley’s integral bound and~\cref{lem:cov} on $\ell_{p}$–covers,
\begin{align*}
\mathfrak{R}(\mathcal{H})
&\le
\frac{12}{\sqrt{n}}
\int_{0}^{F_{\max}}
  \sqrt{\ln\mathcal{N}\left(\mathcal{H},\varepsilon,\|\cdot\|_{1}\right)}
\,\mathrm{d}\varepsilon\\
&\le
\frac{12}{\sqrt{n}}
\int_{0}^{F_{\max}M}
  \sqrt{\ln\mathcal{N}\left(\mathcal{V}_{p,R},\varepsilon,\|\cdot\|_{p}\right)}
\,\mathrm{d}\varepsilon\\
&=
12
\int_{0}^{F_{\max}M}
  \sqrt{\,\frac{d_{\mathcal{V}}}{n}\;\ln\left(\frac{2R\,d_{\mathcal{V}}\,d}{\varepsilon}\right)}
\,\mathrm{d}\varepsilon
+F_{\max}M\,\sqrt{\frac{\ln(2\beta)}{n}}.
\end{align*}

Define
\begin{align*}
I
&=
\int_{0}^{F_{\max}M}
  \sqrt{\ln\left(\frac{2R\,d_{\mathcal{V}}\,d}{\varepsilon}\right)}
\,\mathrm{d}\varepsilon.
\end{align*}

Similar calculation in~\cref{lem:sup_z} shows
\begin{align*}
I
&\le
F_{\max}M
\sqrt{\ln\left(\frac{2R\,d_{\mathcal{V}}\,d}{M\,F_{\max}}\right)}.
\end{align*}
Combining these estimates yields the stated bound.
\end{proof}

\begin{theorem}[Generalization Bound]
Assume that assumptions~\ref{assump:bounded_target} and~\ref{assump:bounded_H} hold (i.e., boundedness of $f_{\bm w} \in \mathcal{H}$, $f^*$, and $\phi$). Let the population and empirical risks be defined as
\begin{align*}
\mathcal{L}(\bm{w}) := \mathbb{E}(f_{\bm w}(X)-Y)^2, \quad
\mathcal{L}_n(\bm{w}) := \frac{1}{n} \sum_{i=1}^n (f_{\bm w}(X_i)-Y_i)^2.
\end{align*}
Then, there exists a universal constant \( C_0, C_1, C_2 > 0 \) such that the following holds with probability at least $1 - \delta$:
\begin{align}
\sup_{\bm{w}} \left| \mathcal{L}_n(\bm{w}) - \mathcal{L}(\bm{w}) \right|
&\leq C_0\left(4F_{\max} + \sigma\sqrt{2/\pi} \right) F_{\max} M 
\left(
\sqrt{\frac{d_{\mathcal{V}}}{n}\log \left(\frac{2Rd_{\mathcal{V}}d}{M F_{\max}} \right)} 
+ \sqrt{\frac{\log(2\beta)}{n}}
\right) \nonumber \\
&\quad + \max \left\{ 
\sqrt{\frac{C_1(4F_{\max}+\sigma)^4}{n} \log \frac{4}{\delta}},\ 
\frac{C_2(4F_{\max}+\sigma)^2 \log n}{n} \log \frac{12}{\delta}
\right\}.
\end{align}
\end{theorem}

\begin{proof}
Apply Adamczak's concentration inequality~\citep{adamczak2008tail} to the supremum:
\begin{align}
\label{eq:adamaczak}
\begin{aligned}
\Pr\Biggl(
\sup_w \left| \frac{1}{n} \sum_{i=1}^n g_w(X_i, Y_i) \right|
&\leq C \mathbb{E} \sup_w \left| \frac{1}{n} \sum_{i=1}^n g_w(X_i, Y_i) \right| + t
\Biggr) \\
&\geq 1 - \left(
\exp\left( - \frac{t^2 n}{\tilde{C}_1 V} \right)
+ \exp\left( - \frac{t n}{\nu}\right)
\right),
\end{aligned}
\end{align}
where $C$ is constant, $V$ and $\nu$ are quantities defined by
\begin{align*}
    V &\coloneqq  \sup_w \mathbb{E}[g_w(X, Y)^2], \\
    \nu &\coloneqq \| \max_{1 \leq i \leq n} \sup_w g_w(X_i, Y_i)\|_{\psi_1}.
\end{align*}

From~\cref{lem:bounded_g,lem:max_subexp}, we have
\begin{gather*}
    \sup_w \mathbb{E}[g_w(X, Y)^2] \leq C_1/\tilde{C}_1 (4F_{\max} + \sigma)^4, \\
    \| \max_{1 \leq i \leq n} \sup_w g_w(X_i, Y_i)\|_{\psi_1} \leq  C_2(4F_{\max} + \sigma)^2 \log n,
\end{gather*}
where $C_1, C_2$ are constants. 

By substituting the above quantities into inequality~\cref{eq:adamaczak} and converting it into a high-probability bound, we obtain, with probability at least $1 - \delta$, the following result:
\begin{align}
\label{eq:adamaczak2}
\sup_w \left| \frac{1}{n} \sum_{i=1}^n g_w(X_i, Y_i) \right|
&\leq C \mathbb{E} \sup_w \left| \frac{1}{n} \sum_{i=1}^n g_w(X_i, Y_i) \right| \nonumber \\
&\quad + \max \left\{
\sqrt{\frac{C_1(4F_{\max} + \sigma)^4}{n} \log \frac{4}{\delta}},\ 
\frac{C_2(4F_{\max} + \sigma)^2 \log n}{n} \log \frac{12}{\delta}
\right\}.
\end{align}

We bound the expectation via symmetrization and Rademacher complexity:
\begin{align}
\label{eq:symmetry_rad}
\mathbb{E} \sup_w \left| \frac{1}{n} \sum_{i=1}^n g_w(X_i, Y_i) \right|
&= \mathbb{E} \sup_w \left| \frac{1}{n} \sum_{i=1}^n 
\left((f(X_i) - Y_i)^2 - (f(X_i') - Y_i')^2 \right) \right| \nonumber \\
&= \mathbb{E} \sup_w \left| \frac{1}{n} \sum_{i=1}^n \tau_i 
\left((f(X_i) - Y_i)^2 - (f(X_i') - Y_i')^2 \right) \right| \nonumber \\
&\leq 2 \mathbb{E} \sup_w \left| \frac{1}{n} \sum_{i=1}^n \tau_i (f(X_i) - Y_i)^2 \right| \nonumber \\
&= 2 \mathfrak{R}(\ell \circ \mathcal{H}),
\end{align}
where $\tau_i$ are Rademacher variables and $\mathfrak{R}(\ell \circ \mathcal{H})$ denotes the Rademacher complexity of the squared loss class.

Combining \cref{lem:contraction_rad,lem:dudley_rad} with \cref{eq:adamaczak2,eq:symmetry_rad}, these bounds complete the proof.
\end{proof}

\section{Generalization Bound for Polynomial PINN}
\label{app:poly_pinn}
In this section, we extend the Rademacher complexity analysis to Physics-Informed Neural Networks (PINNs) constructed with polynomial activation functions.

\begin{definition}[Matrix p-Norm]
A matrix $W \in \mathbb{R}^{d_{\mathrm{in}} \times d_{\mathrm{out}}}$ is treated as a vector $\operatorname{vec}(W) \in \mathbb{R}^{d_{\mathrm{in}}d_{\mathrm{out}}}$, and its norm is measured by the standard vector p-norm:
$$
\|W\|_{p} := \left(\sum_{i=1}^{m}\sum_{j=1}^{n}|W_{ij}|^{p}\right)^{1/p} \quad (1 \le p < \infty).
$$
\end{definition}

\begin{definition}[Mixed Norm for Vector-Valued Functions]
For a function $f: \mathcal{X} \to \mathbb{R}^{d_{\mathrm{out}}}$, its $(\infty, p)$-norm is defined as the supremum of the p-norm of its output vectors over the domain $\mathcal{X}$:
$$
\|f\|_{\infty,p} := \sup_{x \in \mathcal{X}} \|f(x)\|_{p}, \quad \text{where} \quad \|f(x)\|_{p} := \left(\sum_{i=1}^{m}|f_i(x)|^{p}\right)^{1/p}.
$$
\end{definition}

\begin{assumption}[Hypothesis Space of Polynomial PINN]
\label{assump:poly_pinn_space}
Let the hypothesis space $\mathcal{H}_L$ consist of functions $f_{\bm{w}}: \mathbb{R}^{m} \to \mathbb{R}$ realized by a fully-connected neural network of depth $L$:
\begin{align*}
    f_{\bm{w}}(x) = W_L \phi( W_{L-1} \cdots \phi(W_1 x) \cdots ),
\end{align*}
where $\phi$ is a polynomial activation function. 
The parameters $\bm{w} = \operatorname{vec}(W_1, \dots, W_L) \in \mathbb{R}^d$ are constrained to a set $\mathcal{V}_R$:
\begin{align*}
    \bm{w} \in \mathcal{V}_R = \mathcal{V}(\mathscr{D}, \mathcal{T}) \cap \mathbb{B}_2(R),
\end{align*}
where $\mathcal{V}(\mathscr{D},\mathcal{T})$ is the affine variety defined by 
\begin{align*}
\label{eq:def-V-dnn}
\mathcal{V}(\mathscr{D},  \mathcal{T}) \coloneqq \left\{ \bm{w} \in \mathbb{R}^d: \langle \mathscr{D}\left[f_{\bm{w}} \right], \psi_k \rangle_{\mu_k} = 0,\ \forall (\psi_k, \mu_k) \in \mathcal{T}\right\},
\end{align*}
and $\mathbb{B}_2(R) = \{ \bm{w} : \|W_{\ell}\|_{2} \le R,\ 1 \le \forall \ell \le L\}$ is a ball constraining the p-norm of each weight matrix.
\end{assumption}

\begin{assumption}[Bounded Hypotheses and Activation Function]
\label{assump:poly_pinn_bounded}
The polynomial activation function $\phi$ is bounded, $\|\phi\|_{\infty, 2} \le M$ for some constant $M > 0$. Similarly, every hypothesis $f_{\bm{w}} \in \mathcal{H}_L$ is uniformly bounded by $\|f_{\bm{w}}\|_{\infty} \le F_{\max}$.
\end{assumption}

\begin{assumption}[Lipschitz Continuity of Activation Function]
\label{assump:poly_pinn_lip}
The polynomial activation function $\phi$ is Lipschitz continuous with a Lipschitz constant $L_{\phi}$, such that for any $z_1, z_2 \in \mathcal{X}$ in the relevant domain:
\begin{align*}
    \|\phi(z_1) - \phi(z_2)\|_2 \le L_{\phi} \|z_1 - z_2\|_2.
\end{align*}
\end{assumption}

\begin{lemma}[Covering Number of Polynomial PINN]
\label{lem:poly_pinn_covering}
Let $f_{\bm{w}}, f_{\bm{w}'} \in \mathcal{H}_L$ be two functions corresponding to parameters $\bm{w}, \bm{w}' \in \mathcal{V}_R$. If $\|W_{\ell} - W'_{\ell}\|_{2} < \varepsilon$ for all $\ell = 1, \dots, L$, then the distance between the functions is bounded by:
\begin{align*}
    \|f_{\bm{w}} - f_{\bm{w}'}\|_{\infty} \le \ell_L \varepsilon,
\end{align*}
where $\ell_L$ is a stability constant dependent on the network architecture:
\begin{align*}
    \ell_L = M \frac{(R L_{\phi})^{L-1} - 1}{R L_{\phi} - 1} + R(R L_{\phi})^{L-1}.
\end{align*}
Consequently, the covering number of the hypothesis space is bounded by that of the parameter space:
\begin{align*}
    \mathcal{N}(\mathcal{H}_L, \ell_L \varepsilon, \|\cdot\|_{\infty}) \le \mathcal{N}(\mathcal{V}_R, \varepsilon, \|\cdot\|_{2}).
\end{align*}
\end{lemma}

\begin{proof}
Let $h_\ell$ and $h'_\ell$ be the outputs of layer $\ell$ for parameters $\bm{w}$ and $\bm{w}'$ respectively. We can establish a recursive inequality:
\begin{align*}
    \|h_\ell - h'_\ell\|_{\infty, 2} &= \|W_\ell \phi(h_{\ell-1}) - W'_\ell \phi(h'_{\ell-1})\|_{\infty, 2} \\
    &\le \|W_\ell(\phi(h_{\ell-1}) - \phi(h'_{\ell-1}))\|_{\infty, 2} + \|(W_\ell - W'_\ell)\phi(h'_{\ell-1})\|_{\infty, 2} \\
    &\le \|W_\ell\|_{2} L_{\phi} \|h_{\ell-1} - h'_{\ell-1}\|_{\infty, 2} + \|W_\ell - W'_\ell\|_{2} \|\phi(h'_{\ell-1})\|_{\infty, 2} \\
    &\le R L_{\phi} \|h_{\ell-1} - h'_{\ell-1}\|_{\infty, 2} + \varepsilon M.
\end{align*}
Solving this recurrence relation for $\|h_L - h'_L\|_{\infty} = \|f_{\bm{w}} - f_{\bm{w}'}\|_{\infty}$ yields the constant $\ell_L$. The relationship between covering numbers follows directly.
\end{proof}

\begin{lemma}[Dudley Integral Bound for Polynomial PINN]
\label{lem:poly_pinn_dudley}
Let $\mathcal{H}_L$ be the hypothesis space defined in \cref{assump:poly_pinn_space}, where the underlying affine variety $\mathcal{V}$ is $(\beta, d_{\mathcal{V}})$-regular. The Rademacher complexity of $\mathcal{H}_L$ is bounded by:
\begin{align*}
\mathfrak{R}(\mathcal{H}_L)
&\le
C\,F_{\max}\,\ell_L
\left(
\sqrt{\,\frac{d_{\mathcal{V}}}{n}\;\ln\left(\frac{2R\,d_{\mathcal{V}}\,d}{\ell_L\,F_{\max}}\right)}
+\sqrt{\frac{\ln(2\beta)}{n}}
\right).
\end{align*}
for some constant $C>0$.
\end{lemma}

\begin{proof}
The proof follows by applying Dudley's integral theorem and using the covering number bound from \cref{lem:poly_pinn_covering} combined with the bound on $\mathcal{N}(\mathcal{V}_R, \varepsilon, \|\cdot\|_{2})$ from~\cref{lem:cov} (from the main paper).
\begin{align*}
    \mathfrak{R}(\mathcal{H}_L) &\le \frac{12}{\sqrt{n}} \int_0^{F_{\mathrm{max}}} \sqrt{\ln \mathcal{N}(\mathcal{H}_L, \varepsilon, \|\cdot\|_{\infty})} \, d\varepsilon \\
    &\le \frac{12}{\sqrt{n}} \int_0^{\ell_L F_{\mathrm{max}}} \sqrt{\ln \mathcal{N}(\mathcal{V}_R, \varepsilon, \|\cdot\|_{2})} \, d\varepsilon.
\end{align*}
Substituting the known bound for $\ln\mathcal{N}(\mathcal{V}_R, \cdot, \cdot)$ and solving the integral yields the result.
\end{proof}

\begin{theorem}[Generalization Bound for Polynomial PINN]
\label{thm:poly_pinn_gen_bound}
Under assumptions \ref{assump:poly_pinn_space}, \ref{assump:poly_pinn_bounded}, \ref{assump:poly_pinn_lip}, and the standard boundedness of the target function and noise, there exist constants $C_0, C_1, C_2 > 0$ such that with probability at least $1-\delta$:
\begin{align*}
    \sup_{\bm{w} \in \mathcal{V}_R} |\mathcal{L}_n(\bm{w}) - \mathcal{L}(\bm{w})| &\le C_0 \left(4F_{\max} + \sigma\sqrt{2/\pi}\right) F_{\max}\,\ell_L
\left(
\sqrt{\,\frac{d_{\mathcal{V}}}{n}\;\ln\left(\frac{2R\,d_{\mathcal{V}}\,d}{\ell_L\,F_{\max}}\right)}
+\sqrt{\frac{\ln(2\beta)}{n}}
\right) \\
    &\quad + \max\left\{ \sqrt{\frac{C_1(4F_{\max}+\sigma)^4}{n} \log\frac{4}{\delta}}, \frac{C_2(4F_{\max}+\sigma)^2 \log n}{n} \log\frac{12}{\delta} \right\}.
\end{align*}
\end{theorem}

\begin{proof}
The proof structure is identical to that of the main theorem in the previous appendix. We start with Adamczak's concentration inequality. The expectation term is bounded via symmetrization, leading to the Rademacher complexity of the loss class, $\mathfrak{R}(\ell \circ \mathcal{H}_L)$. We apply the contraction principle (\cref{lem:contraction_rad}) to relate this to $\mathfrak{R}(\mathcal{H}_L)$. Finally, we substitute the Dudley integral bound for Polynomial PINNs from \cref{lem:poly_pinn_dudley} to obtain the final generalization bound.
\end{proof}

\begin{remark}[Limitations of~\cref{thm:poly_pinn_gen_bound}]
While~\Cref{thm:poly_pinn_gen_bound} provides a theoretical generalization bound for
polynomial PINNs with neural network surrogates, two caveats are worth
emphasizing.  
First, the affine variety induced by the hard constraints is defined by
a system of high-degree polynomial equations.  The resulting algebraic
structure is computationally intractable to characterize explicitly,
which limits both theoretical validation and practical implementation.  
Second, recent findings suggest that the generalization behavior of
over-parameterized models such as neural networks is not governed solely
by the geometry of the hypothesis space, but is strongly affected by the
implicit bias of the optimization algorithm.  Hence, bounds of the form
in~\cref{thm:poly_pinn_gen_bound} may substantially deviate from the empirically observed performance.
\end{remark}

\section{Experimental Detail}
\label{app:exp_detail}

\subsection{Experiments on Strong Solution}
\label{app:exp-strong-sol}
In the experiments in \cref{sec:strong-sol}, strong solutions to the equations are obtained analytically. The analytical solution with added Gaussian noise was used as data, the variance of the Gaussian noise was set to \(0.01\). The hyperparameters \(L^2\) regularization weights and differential equation constraint weights \(\xi\) and \(\nu\) were searched in the range \([1\text{e-}9, 1\text{e-}2]\) using the Optuna library \citep{akiba2019optuna}. The configuration with the smallest MSE on the validation data among 100 candidates was selected. All experiments were conducted on a MacBook Air equipped with an Apple M3 chip and 64 GB of unified memory. No external GPU or cluster computing resources were used.


\textbf{Harmonic Oscillator:} The initial value problem of a harmonic oscillator \(\mathscr{D}[y]=0\) with spring constant \( k_s \) and mass \( m_s \) on the domain \( \Omega = [0, T] \) is given by:
\begin{align*}
    \mathscr{D}[y] = \frac{\mathrm{d}^2}{\mathrm{d}t^2}y + \frac{k_s}{m_s}y,\quad y(0)=y_0,\quad \frac{\mathrm{d}}{\mathrm{d}t}y(0)=v_0.
\end{align*}
We set the parameters \(m_s = k_s = 1.0,\ T = 2\pi\). The initial position and velocity \([y_0, v_0]^{\top}\) are generated from the normal distribution \(\mathcal{N}(\bm{1}, I)\), where \(\bm{1}\) is an all-ones vector and \(I\) is the identity matrix.
The solution to the initial value problem is analytically given by:
\begin{align*}
    y(t) = y_0 \cos(\omega t) + \frac{v_0}{\omega} \sin(\omega t), \ \omega = \sqrt{k_s/m_s}.
\end{align*}
The settings for the basis functions and the trial functions with the measure \(\phi_j \in \mathcal{B},\ (\psi_k, \mu_k) \in \mathcal{T}\) are as follows:
\begin{gather*}
    \phi_1(x) = 1,\
    \phi_{2j}(x) = \cos\left(\frac{2 \pi j}{T}x\right),\ \phi_{2j+1}(x) = \sin\left(\frac{2 \pi j}{T}x\right) \ (j = 1,\ \ldots,\ d_t),\\
    \psi_k(x) = 1,\ \mu_k = \delta_{x_k} \ (k = 1,\ \ldots,\ K),
\end{gather*}
where \(d_t \in \{2, 4, 8, 16\}\) is the set of the number of basis functions, and \(x_k \in \Omega\) is uniformly sampled from data with \(K = 100\).

\textbf{Diffusion Equation:} The initial value problem for the one-dimensional diffusion equation \(\mathscr{D}[u]=0\) with diffusion coefficient \( c \) and periodic boundary conditions is given by:
\begin{gather*}
\begin{aligned}
\mathscr{D}[u] = \frac{\partial}{\partial t} u - c\frac{\partial^2}{\partial x^2} u & \ & (x, t) \in [-\Xi, \Xi] \times [0, T],
u(x, 0)=u_0(x) & \quad & x \in [- \Xi, \Xi]
\end{aligned} \\
u\left(- \Xi,t\right)=u\left( \Xi,t\right),\quad \frac{\partial u}{\partial x}\left(- \Xi, t\right)= \frac{\partial u}{\partial x}\left( \Xi, t\right).
\end{gather*}

We set the parameters \(c = 1.0,\ \Xi = \pi,\ T=2\pi\). The initial value \(u_0\) is given by:
\begin{align}
    \label{eq:diffusion-init}
    u_0(x) = \sum_{j=0}^{j_{\mathrm{max}}} A_j \cos\left(\omega_j x\right)+B_j \sin\left(\omega_j x\right),\ \omega_j = \frac{j\pi}{\Xi},
\end{align}
where \([A_j, B_j]^{\top}\) are generated from the normal distribution \(\mathcal{N}(\bm{1}, I)\) for all \(j = 0,\ \ldots, j_{\mathrm{max}}\) and \(j_{\mathrm{max}}\) is set to \(1\).
The solution to the initial value problem is analytically given by:
\begin{align*}
    u(x, t) = \sum_{j=0}^{j_{\mathrm{max}}}\left[A_j \cos\left(\omega_j x\right)+B_j \sin\left(\omega_j x\right)\right] e^{-c \omega_j^2 t}.
\end{align*}
The settings for the basis functions and the trial functions with the measure \(\phi_j \in \mathcal{B},\ (\psi_k, \mu_k) \in \mathcal{T}\) are as follows:
\begin{gather*}
    \phi_1(x, t) = 1,\
    \phi_{2j, j'}(x, t) = \cos(\omega_j x) e^{-c\omega_{j'}^2 t},\ \phi_{2j+1, j'}(x, t) = \sin(\omega_j x) e^{-c \omega_{j'}^2 t} \\
    (j=1,\ \ldots,\ d_x,\ j' = 1,\ \ldots,\ d_t),\\
    \psi_k(x, t) = 1,\ \mu_k = \delta_{(x_k, t_k)} \ (k = 1,\ \ldots,\ K),
\end{gather*}
where \(d_t = 2,\ d_x \in \{10, 15, 20, 25\}\) are the sets of the number of basis functions, and \((x_k, t_k) \in \Omega\) is uniformly sampled from data with \(K = 50 \times 500\).

\subsection{Experiments on Numerical Solution}
\label{app:exp-numel-sol}
In the experiments in~\cref{sec:numel-sol}, we numerically simulate the Bernoulli equation using the explicit Euler method and the diffusion equation using the finite difference method (FDM). The data used are the numerical solutions with added Gaussian noise of variance 0.01. The method for hyperparameter search is the same as described in~\cref{app:exp-strong-sol}.
For the nonlinear equations, we use the Adam optimizer with a learning rate of \(1 \times 10^{-2}\), along with an exponential learning rate scheduler. The training is performed for a maximum of 2000 epochs, utilizing an early stopping technique.

\textbf{Discrete Bernoulli Equation:}
The discrete Bernoulli equation \(\mathscr{D}_{h}[y]=0\) with the step size \( h \) on the domain \(\Omega = [0, T]\) is given by:
\begin{align*}
    \mathscr{D}_h[y] &= \frac{y_{\tau+1} - y_{\tau}}{h} + Py_{\tau} - Qy_{\tau}^{\rho},
\end{align*}
where \( y_{\tau} = y(t_{\tau}) \) and \( y_{\tau+1} = y(t_{\tau} + h) \) are evaluations on the grid \( \{ t_{\tau} \}_{\tau=1}^{n_t} \) with \( n_t = \tfrac{T}{h} \). 
We set the constant parameters \( (P, Q, \rho) \) to \( (1.0, 0.0, 0.0) \) for the linear case and to \( (1.0, 0.5, 2.0) \) for the non-linear case. We use varying \( n_t \in \{ 100, 200 \} \) with \( T = 1.0 \) for both cases. The initial state \( y_0 \) is generated from the standard normal distribution \(\mathcal{N}(0, 1)\) for both cases.
The ground-truth solution to the initial value problem is numerically solved by the explicit Euler method with step size \( h \).
The settings for the basis functions and the trial functions with measure \(\phi_{\tau} \in \mathcal{B}_h\), \((\psi_{\tau}, \mu_{\tau}) \in \mathcal{T}_h\) are as follows:
\begin{gather*}
    \phi_{\tau}(t) = \begin{cases}
    1 & \text{if } t \in [t_{\tau}, t_{\tau+1}) \\
    0 & \text{otherwise}
    \end{cases} \quad (\tau = 1, \ldots, n_t), \\
    \psi_{\tau}(t) = \phi_{\tau}(t), \quad \mu_{\tau} = \delta_{t_{\tau}} \quad (\tau = 1, \ldots, n_t),
\end{gather*}
where \(n_t = \frac{T}{h} \) is the same as the number of basis and trial functions, corresponding to the ground-truth solutions.

\textbf{Discrete Diffusion Equation:}
The one-dimensional discrete diffusion equation \(\mathscr{D}_{\bm{h}}[u]=0\) with the step size \(\bm{h} = [h_t, h_x]^{\top}\) and the diffusion coefficient \(c(u)\) on the domain \(\Omega = [-\Xi, \Xi] \times [0, T]\) is given by:
\begin{align*}
    \mathscr{D}_{\bm{h}}[u] = \frac{u^{\tau+1}_{j} - u^{\tau}_{j}}{h_t} - c(u^{\tau}_{j}) \frac{u^{\tau}_{j+1} - 2u^{\tau}_{j} + u^{\tau}_{j-1}}{h_x^2},
\end{align*}
where \(u^{\tau}_j \coloneqq u(x_j, t_{\tau})\), \(u^{\tau+1}_j \coloneqq u(x_j, t_{\tau} + h_t)\), and \(u^{\tau}_{j\pm 1} \coloneqq u(x_j \pm h_x, t_{\tau})\) are evaluations on the \(n_x \times n_t\) size grid \( \{ x_j \}_{j=1}^{n_x} \times \{ t_{\tau} \}_{\tau=1}^{n_t} \), where \(n_x \coloneqq \frac{2\Xi}{h_x}\) and \(n_t \coloneqq \frac{T}{h_t}\). The periodic boundary condition is adopted in the spatial domain, \ie, \(u_{n_x + j}^{\tau} = u_{j}^{\tau}\) for any \(j \in [d]\).
The diffusion coefficient \(c(u) = 1.0\) is used for the linear case and \(c(u) = 0.1/(1 + u^2)\) for the nonlinear case. We use varying \( (n_t, n_x) \in \{ (400, 10), (400, 20), (400, 30) \}\) with \(\Xi = 1.0\) and \(T = 1.0\) for both cases. The initial value is generated with the same setting as shown in \cref{eq:diffusion-init}.
The ground-truth solution to the initial value problem is numerically solved by the FDM with step sizes \(h_t\) for the time domain and \(h_x\) for the spatial domain.
The settings for the basis functions and the trial functions with measure \(\phi_{j, \tau} \in \mathcal{B}_{\bm{h}},\ (\psi_{j, \tau}, \mu_{j, \tau}) \in \mathcal{T}_{\bm{h}}\) are as follows:
\begin{gather*}
    \phi_{j, \tau}(x, t) = \begin{cases}
    1 & \text{if } (x, t) \in [x_j, x_{j+1}] \times [t_{\tau}, t_{\tau+1}] \\
    0 & \text{otherwise}
    \end{cases} \quad (j = 1, \ldots, n_x, \ \tau = 1, \ldots, n_t), \\
    \psi_{j, \tau}(x, t) = \phi_{j, \tau}(x, t), \quad \mu_{j, \tau} = \delta_{(x_j, t_{\tau})} \quad (j = 1, \ldots, n_x, \ \tau = 1, \ldots, n_t),
\end{gather*}
where \(n_x = \frac{2\Xi}{h_x}\) and \(n_t = \frac{T}{h_t}\) are the same as the number of basis and trial functions, corresponding to the ground-truth solutions.

\section{Additional Experimental Results}
\label{app:additional_exp}
For each benchmark (discrete linear/nonlinear Bernoulli and Heat equations), we fix the number of basis functions \(d\) and vary the size of the trial‐function set \(\mathcal{T}\).  Reducing \(|\mathcal{T}|\) relaxes the algebraic constraints on the learned solution, which in turn increases the dimension of the associated affine variety \(d_{\mathcal{V}}\). As reported in Tables~\cref{fig:bernoulli-pilr-sub,fig:heat-pilr-sub}, when \(|\mathcal{T}|\) decreases (and thus \(d_{\mathcal{V}}\) increases), the Test MSE steadily increases.  This consistent rising trend of Test MSE with larger \(d_{\mathcal{V}}\) demonstrates that models endowed with fewer trial functions (i.e.\ weaker constraints) generalize more poorly.
% --- Bernoulli の結果 ---
\begin{figure*}[!htb]
  \centering
  \captionsetup[subfigure]{justification=centering}
  % Linear
  \begin{subfigure}[t]{\textwidth}
    \centering
    \caption{Linear Bernoulli eq.}
    \begin{tabular}{cc|ccc}
      \toprule
      Settings & $h$ & \multicolumn{3}{c}{$1/100$} \\
      \midrule
      \multirow{2}{*}{Dimensions}  & $d$                & \multicolumn{3}{c}{100}       \\
                                   & $d_{\mathcal{V}}$  & 10 & 20 & 40 \\
      \midrule
      Test MSE (PILR)             &                    & $0.012 \pm 0.0023$ & $0.13 \pm 0.082$ & $0.33 \pm 0.22$ \\
      \bottomrule
    \end{tabular}
  \end{subfigure}

  \vspace{1em}

  % Nonlinear
  \begin{subfigure}[t]{\textwidth}
    \centering
    \caption{Nonlinear Bernoulli eq.}
    \begin{tabular}{cc|ccc}
      \toprule
      Settings & $h$ & \multicolumn{3}{c}{$1/100$} \\
      \midrule
      \multirow{2}{*}{Dimensions}  & $d$                & \multicolumn{3}{c}{100}       \\
                                   & $d_{\mathcal{V}}$  & 10 & 20 & 40 \\
      \midrule
      Test MSE (PILR)             &                    & $0.17 \pm 0.11$ & $0.21 \pm 0.14$ & $0.33 \pm 0.23$ \\
      \bottomrule
    \end{tabular}
  \end{subfigure}

  \caption{Experimental results for PILR on the discrete Bernoulli equations.}
  \label{fig:bernoulli-pilr-sub}
\end{figure*}


% --- Heat の結果 ---
\begin{figure*}[!htb]
  \centering
  \captionsetup[subfigure]{justification=centering}
  % Linear
  \begin{subfigure}[t]{\textwidth}
    \centering
    \caption{Linear Heat eq.}
    \begin{tabular}{cc|ccc}
      \toprule
      Settings & $(h_t,h_x)$ & \multicolumn{3}{c}{$(1/400,\,2/10)$} \\
      \midrule
      \multirow{2}{*}{Dimensions}  & $d$                & \multicolumn{3}{c}{4010}      \\
                                   & $d_{\mathcal{V}}$  & 110 & 210 & 410 \\
      \midrule
      Test MSE (PILR)             &                    & $1.6 \pm 0.35$ & $1.9 \pm 0.44$ & $2.0 \pm 0.49$ \\
      \bottomrule
    \end{tabular}
  \end{subfigure}

  \vspace{1em}

  % Nonlinear
  \begin{subfigure}[t]{\textwidth}
    \centering
    \caption{Nonlinear Heat eq.}
    \begin{tabular}{cc|ccc}
      \toprule
      Settings & $(h_t,h_x)$ & \multicolumn{3}{c}{$(1/200,\, 2/10)$} \\
      \midrule
      \multirow{2}{*}{Dimensions}  & $d$                & \multicolumn{3}{c}{2010}      \\
                                   & $d_{\mathcal{V}}$  & 110 & 210 & 410 \\
      \midrule
      Test MSE (PILR)             &                    & $0.37 \pm 0.11$ & $0.43 \pm 0.14$ & $0.56 \pm 0.19$ \\
      \bottomrule
    \end{tabular}
  \end{subfigure}

  \caption{Experimental results for PILR on the discrete Heat equations.}
  \label{fig:heat-pilr-sub}
\end{figure*}
\end{document}
